% !TeX program = lualatex
% ================================================================
% Urdu Parallel Text Version of FactorsMultiplesPrimesQuickQuestions
%
% Generated by the server using the following settings:
%
%   language            Urdu (urd)
%   text direction      right to left
%   font                Noto Nastaliq Urdu
%   fallback font       Noto Serif
%   built from          latex/quickQuestions/factorsMultiplesPrimesQuickQuestions.tex
%   vocabulary key      latex/quickQuestions/factorsMultiplesPrimesQuickQuestions/L2/urd/factorsMultiplesPrimesQuickQuestions_urduKey.tex
%   tier 3 vocabulary   green   RGB 0 128 0
%   English text        black   RGB 0 0 0
%   Urdu text           purple  RGB 112 48 160
%   layout macros       version 3
%
% Please don't edit any information in this comment block - the server
% compares it against the current settings and rebuilds this file when
% they differ, so changes here would be overwritten. However, feel free
% to make updates to the LaTeX below if you want to edit it.
% ================================================================

\documentclass[aspectratio=169]{beamer}
% ================================================================
% Copied from the English sheet this was translated from, so that
% anything it sets up for itself still works here
% ================================================================
\usepackage{amsmath}
\usepackage{amssymb}
\usepackage{tikz}
\usetikzlibrary{shapes.geometric}

% ================================================================
% Quick Questions deck: factors, multiples and primes.
%
% Each topic opens with five true/false statements and then runs ten
% quick questions that get gradually harder. Every question gets two
% slides: the question on its own for the mini whiteboards, then the
% same question again with the solution underneath in red.
%
% Two topics are built differently. "Is it prime?" runs every whole
% number from 1 to 100 on its own slide, then lists that number's
% factors and says why it is or is not prime. "HCF and LCM from lists"
% builds each answer over five slides, adding the factor lists, the
% multiples and then the rings round the HCF and the LCM one slide at
% a time.
% ================================================================

\setbeamertemplate{navigation symbols}{}

% A link in the bottom left of every slide back to the contents, so a
% deck this long can be jumped around in during a lesson.
\setbeamertemplate{footline}{%
  \hbox to\paperwidth{%
    \hspace{1em}%
    \hyperlink{contents}{%
      \color{gray}\scriptsize$\leftarrow$~Back to contents}%
    \hfil}%
  \vskip4pt%
}

% Topic divider.
\newcommand{\qtopic}[1]{%
  \section{#1}%
  \begin{frame}
    \vfill
    \begin{center}
      {\usebeamercolor[fg]{structure}\Huge\bfseries #1}
    \end{center}
    \vfill
  \end{frame}%
}

% \qq[<question size>][<solution size>]{<question>}{<solution>}
% The two optional sizes drop the type for the wordier questions and
% the longer solutions, so nothing runs off the slide.
\NewDocumentCommand{\qq}{O{\Huge}O{\LARGE}+m+m}{%
  \begin{frame}
    \vfill
    \begin{center}
      \begin{minipage}{0.92\textwidth}
        \centering #1 #3
      \end{minipage}
    \end{center}
    \vfill
  \end{frame}%
  \begin{frame}
    \vfill
    \begin{center}
      \begin{minipage}{0.92\textwidth}
        \centering #1 #3
        \par\vspace{0.9em}
        {\color{red}#2 #4}
      \end{minipage}
    \end{center}
    \vfill
  \end{frame}%
}

% \tf[<statement size>][<answer size>]{<statement>}{<answer>}
% As \qq, but with a standing "True or false?" heading above the
% statement. The statements hunt for the usual misconceptions, so the
% answer always says why, not just which.
\NewDocumentCommand{\tf}{O{\LARGE}O{\Large}+m+m}{%
  \begin{frame}
    \vfill
    \begin{center}
      \begin{minipage}{0.92\textwidth}
        \centering
        {\usebeamercolor[fg]{structure}\Large\bfseries True or false?}
        \par\vspace{1em}
        #1 #3
      \end{minipage}
    \end{center}
    \vfill
  \end{frame}%
  \begin{frame}
    \vfill
    \begin{center}
      \begin{minipage}{0.92\textwidth}
        \centering
        {\usebeamercolor[fg]{structure}\Large\bfseries True or false?}
        \par\vspace{1em}
        #1 #3
        \par\vspace{0.9em}
        {\color{red}#2 #4}
      \end{minipage}
    \end{center}
    \vfill
  \end{frame}%
}

% \prq{<number>}{<factors>}{<why it is or is not prime>}
% The "Is it prime?" drill: the number alone in a large centred font,
% then the same number with all of its factors and the reason in red.
% The number is scaled rather than set with \fontsize, because the
% sans font stops at 51.59pt and larger sizes are silently substituted.
\NewDocumentCommand{\prq}{m m +m}{%
  \begin{frame}
    \vfill
    \begin{center}
      \scalebox{3.6}{\Huge\bfseries #1}
    \end{center}
    \vfill
  \end{frame}%
  \begin{frame}
    \vfill
    \begin{center}
      \scalebox{1.9}{\Huge\bfseries #1}
      \par\vspace{0.8em}
      \begin{minipage}{0.9\textwidth}
        \centering\color{red}
        {\large Factors of $#1$: #2}
        \par\vspace{0.7em}
        {\large #3}
      \end{minipage}
    \end{center}
    \vfill
  \end{frame}%
}

% Rings drawn round a number in a list. The ring is set in a box the
% width of the number itself, so the list is spaced as though it were
% ordinary text and nothing moves on the slide where the ring appears;
% the ring itself spills over the gaps either side, as a ring drawn by
% hand would.
\tikzset{
  ringed/.style={ellipse, inner sep=1pt, draw=red, line width=0.9pt},
  unringed/.style={ellipse, inner sep=1pt, draw=none},
}
\newlength{\ringwd}
\newcommand{\ringnode}[2]{%
  \settowidth{\ringwd}{$#2$}%
  \addtolength{\ringwd}{3pt}% a little clearance for the ring to sit in
  \makebox[\ringwd][c]{%
    \tikz[baseline=(n.base)]{\node[#1] (n) {$#2$};}%
  }%
}
\newcommand{\hcfmark}[1]{\alt<4->{\ringnode{ringed}{#1}}{\ringnode{unringed}{#1}}}
\newcommand{\lcmmark}[1]{\alt<5->{\ringnode{ringed}{#1}}{\ringnode{unringed}{#1}}}

% \hcflcm{<a>}{<b>}{<factors of a>}{<factors of b>}
%        {<multiples of a>}{<multiples of b>}
% Five slides: the question, the two factor lists, the two lists of
% multiples, the ring round the HCF, then the ring round the LCM.
% Everything added after the question slide is in red.
% The frame is top aligned so that the question stays put as the lists
% are added underneath it.
\NewDocumentCommand{\hcflcm}{m m +m +m +m +m}{%
  \begin{frame}<1-5>[t]
    \vspace{0.8em}
    \begin{center}
      {\LARGE Find the HCF and LCM of $#1$ and $#2$.}
    \end{center}
    \vspace{1.6em}
    \begin{columns}[T]
      \begin{column}{0.40\textwidth}
        \only<2->{\color{red}\large
          Factors of $#1$:\\[0.35em]
          #3
          \par\vspace{1em}
          Factors of $#2$:\\[0.35em]
          #4}
      \end{column}
      \begin{column}{0.58\textwidth}
        \only<3->{\color{red}\large
          Multiples of $#1$:\\[0.35em]
          #5
          \par\vspace{1em}
          Multiples of $#2$:\\[0.35em]
          #6}
      \end{column}
    \end{columns}
    \vfill
  \end{frame}%
}

% Contents-slide bullets. Green then orange sets the difficulty band;
% blue marks the two topics that work through a method a slide at a
% time rather than asking for a quick answer.
\newcommand{\tocbullet}[1]{\textcolor{#1}{$\bullet$}}
\setbeamertemplate{section in toc}{%
  \ifcase\inserttocsectionnumber
    \or\tocbullet{green!55!black}% 1 Factors
    \or\tocbullet{green!55!black}% 2 Multiples
    \or\tocbullet{green!55!black}% 3 Primes
    \or\tocbullet{blue}%          4 Is it prime?
    \or\tocbullet{orange}%        5 Prime factor decomposition
    \or\tocbullet{blue}%          6 HCF and LCM from lists
    \or\tocbullet{orange}%        7 HCF and LCM by prime factor decomposition
    \or\tocbullet{orange}%        8 Square numbers from prime factors
    \or\tocbullet{orange}%        9 Problem solving with HCF and LCM
  \fi
  ~\inserttocsection%
}

\title{Factors, Multiples and Primes: Quick Questions}
\subtitle{Mini whiteboard questions, with solutions}
\author{}
\date{}

% Characters this language's font has not got are borrowed from another,
% rather than being silently dropped - see docs/translated-sheet-latex.md
\directlua{%
  if luaotfload and luaotfload.add_fallback then
    luaotfload.add_fallback("eallatinfallback",
      { "Noto Serif:mode=harf;" })
  else
    tex.error("this luaotfload is too old to register a fallback font")
  end
}
\usepackage[english,bidi=basic]{babel}
\babelprovide[import]{urdu}
\babelfont[urdu]{rm}[Renderer=Harfbuzz, BoldFont={Noto Nastaliq Urdu}, ItalicFont={Noto Nastaliq Urdu}, BoldItalicFont={Noto Nastaliq Urdu}, RawFeature={fallback=eallatinfallback}]{Noto Nastaliq Urdu}
\babelfont[urdu]{sf}[Renderer=Harfbuzz, BoldFont={Noto Nastaliq Urdu}, ItalicFont={Noto Nastaliq Urdu}, BoldItalicFont={Noto Nastaliq Urdu}, RawFeature={fallback=eallatinfallback}]{Noto Nastaliq Urdu}
\newcommand{\ealtext}[1]{\foreignlanguage{urdu}{#1}}
\newcommand{\ealtextblock}[1]{{\raggedleft\foreignlanguage{urdu}{#1}\par}}
% ================================================================
% EAL layout helpers
% ================================================================
\usepackage{xcolor}

\definecolor{ealkeycolour}{RGB}{0,128,0}
\definecolor{ealenglish}{RGB}{0,0,0}
\definecolor{ealtwocolour}{RGB}{112,48,160}

\newsavebox{\ealboxen}
\newsavebox{\ealboxtr}
\newlength{\ealglosswd}
\newlength{\ealglossraise}
\setlength{\ealglossraise}{1.15em}

% a tier 3 word, in English and in translation alike
\newcommand{\ealkey}[1]{\textcolor{ealkeycolour}{#1}}
\newcommand{\ealkeytr}[1]{\textcolor{ealkeycolour}{\ealtext{#1}}}

% One translated word above one English word. Built from kernel
% primitives, never a tabular, which would break wherever the sheet
% loads array, booktabs or tabularx. The box is as wide as the wider of
% the two, so a long translation pushes its neighbours aside instead of
% overlapping them, and it claims its full height so a table row or a
% TikZ node makes room for it.
\newcommand{\ealgloss}[2]{%
  \sbox{\ealboxen}{\ealkey{#1}}%
  \sbox{\ealboxtr}{\small\textcolor{ealtwocolour}{\ealtext{#2}}}%
  \setlength{\ealglosswd}{\wd\ealboxen}%
  \ifdim\wd\ealboxtr>\ealglosswd\setlength{\ealglosswd}{\wd\ealboxtr}\fi
  \leavevmode
  \makebox[\ealglosswd][c]{%
    \makebox[0pt][c]{\usebox{\ealboxen}}%
    \makebox[0pt][c]{\raisebox{\ealglossraise}{\usebox{\ealboxtr}}}%
  }%
}

% opens up line spacing around a block of glosses, so the raised
% translations sit clear of the line above
\newenvironment{ealglossed}{\par\linespread{2.1}\selectfont\ignorespaces}{\par}

% a whole translated sentence above its English counterpart. block level,
% so it does not belong inside a tabular cell
\newcommand{\ealpara}[2]{%
  \par\addvspace{0.5\baselineskip}%
  {\color{ealtwocolour}\ealtextblock{#1}}%
  \penalty200
  {\color{ealenglish}#2\par}%
  \addvspace{0.5\baselineskip}%
}

\begin{document}

\frame{\titlepage}

\begin{frame}[label=contents]{Contents}
  \small
  \tableofcontents
\end{frame}

%================================================================
\qtopic{Factors}
%================================================================

\tf{\ealpara{$7$، $21$ کا \ealkeytr{عامل} ہے}{$7$ is a \ealkey{factor} of $21$}}
  {\ealpara{\textbf{سچ۔} $21=7\times 3$، اس لیے $21$ کو $7$ سے پورا تقسیم کیا جا سکتا ہے۔}{\textbf{True.} $21=7\times 3$, so $21$ divides exactly by $7$.}}

\tf{\ealpara{$1$ ہر مکمل عدد کا \ealkeytr{عامل} ہے}{$1$ is a \ealkey{factor} of every whole number}}
  {\ealpara{\textbf{سچ۔} ہر عدد $n$ کو $1\times n$ لکھا جا سکتا ہے۔}{\textbf{True.} Every number $n$ can be written as $1\times n$.}}

\tf{\ealpara{$6$، $16$ کا \ealkeytr{عامل} ہے}{$6$ is a \ealkey{factor} of $16$}}
  {\ealpara{\textbf{جھوٹ۔} $16$ کے \ealkeytr{عوامل} $1,2,4,8,16$ ہیں؛ $16\div 6$ مکمل عدد نہیں ہے۔}{\textbf{False.} The \ealkey{factors} of $16$ are $1,2,4,8,16$;\\[0.3em] $16\div 6$ is not a whole number.}}

\tf{\ealpara{ہر عدد کے \ealkeytr{عوامل} کی تعداد جفت ہوتی ہے}{Every number has an even number of \ealkey{factors}}}
  {\ealpara{\textbf{جھوٹ۔} \ealkeytr{عوامل} عام طور پر جوڑے بناتے ہیں، لیکن $9$ کے پاس صرف $1,3,9$ ہیں --- تین \ealkeytr{عوامل} --- کیونکہ $3\times 3$ ایک ہی \ealkeytr{عامل} کو دو بار استعمال کرتا ہے۔}{\textbf{False.} \ealkey{Factors} usually pair up, but $9$ has just $1,3,9$ --- three \ealkey{factors} --- because $3\times 3$ uses the same \ealkey{factor} twice.}}

\tf{\ealpara{اگر $4$ کسی عدد کا \ealkeytr{عامل} ہو، تو $8$ بھی اس کا \ealkeytr{عامل} ہوگا}{If $4$ is a \ealkey{factor} of a number,\\[0.3em]then $8$ must be a \ealkey{factor} of it too}}
  {\ealpara{\textbf{جھوٹ۔} $4$، $12$ کا \ealkeytr{عامل} ہے، لیکن $8$ نہیں ہے۔}{\textbf{False.} $4$ is a \ealkey{factor} of $12$, but $8$ is not.}}

\qq[\LARGE]{\ealpara{$12$ کے تمام \ealkeytr{عوامل} لکھیں}{List all the \ealkey{factors} of $12$}}{$1,\ 2,\ 3,\ 4,\ 6,\ 12$}

\qq[\LARGE]{\ealpara{$20$ کے تمام \ealkeytr{عوامل} لکھیں}{List all the \ealkey{factors} of $20$}}{$1,\ 2,\ 4,\ 5,\ 10,\ 20$}

\qq[\LARGE]{\ealpara{$36$ کے تمام \ealkeytr{عوامل} لکھیں}{List all the \ealkey{factors} of $36$}}
  {$1,\ 2,\ 3,\ 4,\ 6,\ 9,\ 12,\ 18,\ 36$}

\qq[\LARGE]{\ealpara{$16$ کے کتنے \ealkeytr{عوامل} ہیں؟}{How many \ealkey{factors} does $16$ have?}}
  {\ealpara{$1,\ 2,\ 4,\ 8,\ 16$\\[0.3em]پانچ \ealkeytr{عوامل}۔}{$1,\ 2,\ 4,\ 8,\ 16$\\[0.3em]Five \ealkey{factors}.}}

\qq[\LARGE]{\ealpara{$12$ اور $18$ کا سب سے بڑا مشترک عامل تلاش کریں}{Find the \ealkey{highest common factor} of $12$ and $18$}}
  {\ealpara{$12$ کے \ealkeytr{عوامل}: $1,\ 2,\ 3,\ 4,\ \mathbf{6},\ 12$\\[0.3em]
   $18$ کے \ealkeytr{عوامل}: $1,\ 2,\ 3,\ \mathbf{6},\ 9,\ 18$\\[0.3em]
   $\mathrm{HCF}=6$}{\ealkey{Factors} of $12$: $1,\ 2,\ 3,\ 4,\ \mathbf{6},\ 12$\\[0.3em]
   \ealkey{Factors} of $18$: $1,\ 2,\ 3,\ \mathbf{6},\ 9,\ 18$\\[0.3em]
   $\mathrm{HCF}=6$}}

\qq[\LARGE]{\ealpara{$24$ اور $36$ کا سب سے بڑا مشترک عامل تلاش کریں}{Find the \ealkey{highest common factor} of $24$ and $36$}}
  {\ealpara{$24$ کے \ealkeytr{عوامل}: $1,\ 2,\ 3,\ 4,\ 6,\ 8,\ \mathbf{12},\ 24$\\[0.3em]
   $36$ کے \ealkeytr{عوامل}: $1,\ 2,\ 3,\ 4,\ 6,\ 9,\ \mathbf{12},\ 18,\ 36$\\[0.3em]
   $\mathrm{HCF}=12$}{\ealkey{Factors} of $24$: $1,\ 2,\ 3,\ 4,\ 6,\ 8,\ \mathbf{12},\ 24$\\[0.3em]
   \ealkey{Factors} of $36$: $1,\ 2,\ 3,\ 4,\ 6,\ 9,\ \mathbf{12},\ 18,\ 36$\\[0.3em]
   $\mathrm{HCF}=12$}}

\qq[\LARGE]{\ealpara{$48$ کے تمام \ealkeytr{عوامل کے جوڑے} لکھیں}{List all the \ealkey{factor pairs} of $48$}}
  {$1\times 48\qquad 2\times 24\qquad 3\times 16$\\[0.4em]
   $4\times 12\qquad 6\times 8$}

\qq[\LARGE]{\ealpara{$30$ سے نیچے کون سا مکمل عدد سب سے زیادہ \ealkeytr{عوامل} رکھتا ہے؟}{Which whole number below $30$ has the most \ealkey{factors}?}}
  {\ealpara{$24$، آٹھ \ealkeytr{عوامل} کے ساتھ: $1,\ 2,\ 3,\ 4,\ 6,\ 8,\ 12,\ 24$}{$24$, with eight \ealkey{factors}: $1,\ 2,\ 3,\ 4,\ 6,\ 8,\ 12,\ 24$}}

\qq[\LARGE]{\ealpara{$n$ کے بالکل تین \ealkeytr{عوامل} ہیں۔ $n$ کی ایک ممکنہ قیمت بتائیں۔}{$n$ has exactly three \ealkey{factors}.\\[0.4em] Give a possible value of $n$.}}
  {\ealpara{$4$ ($1,2,4$)، $9$ ($1,3,9$) یا $25$ ($1,5,25$)۔ کوئی بھی \ealkeytr{مفرد عدد} کا مربع کام کرتا ہے۔}{$4$ ($1,2,4$), $9$ ($1,3,9$) or $25$ ($1,5,25$).\\[0.3em] Any \ealkey{prime} squared works.}}

\qq[\Large][\Large]{\ealpara{وضاحت کریں کہ \ealkeytr{مربع اعداد} کے \ealkeytr{عوامل} کی تعداد طاق کیوں ہوتی ہے}{Explain why \ealkey{square numbers} have an odd number of \ealkey{factors}}}
  {\ealpara{\ealkeytr{عوامل} جوڑوں میں آتے ہیں، ایک \ealkeytr{جذر} کے دونوں طرف۔ \ealkeytr{مربع عدد} میں \ealkeytr{جذر} خود سے جوڑا بناتا ہے، اس لیے ایک \ealkeytr{عامل} بغیر جوڑے کے رہ جاتا ہے۔}{\ealkey{Factors} come in pairs, one either side of the \ealkey{square root}.\\[0.3em] In a \ealkey{square number} the \ealkey{square root} pairs with itself, so one \ealkey{factor} is left unpaired.}}

%================================================================
\qtopic{Multiples}
%================================================================

\tf{\ealpara{$24$، $6$ کا \ealkeytr{مضاعف} ہے}{$24$ is a \ealkey{multiple} of $6$}}
  {\ealpara{\textbf{سچ۔} $6\times 4=24$۔}{\textbf{True.} $6\times 4=24$.}}

\tf{\ealpara{$10$ کا ہر \ealkeytr{مضاعف} $5$ کا بھی \ealkeytr{مضاعف} ہے}{Every \ealkey{multiple} of $10$ is also a \ealkey{multiple} of $5$}}
  {\ealpara{\textbf{سچ۔} $10=5\times 2$، اس لیے $10$ کے ہر \ealkeytr{مضاعف} میں پہلے سے $5$ کا \ealkeytr{عامل} موجود ہے۔}{\textbf{True.} $10=5\times 2$, so every \ealkey{multiple} of $10$ already contains a \ealkey{factor} of $5$.}}

\tf{\ealpara{$35$، $4$ کا \ealkeytr{مضاعف} ہے}{$35$ is a \ealkey{multiple} of $4$}}
  {\ealpara{\textbf{جھوٹ۔} $4\times 8=32$ اور $4\times 9=36$، اس لیے $35$ چھوڑ دیا گیا ہے۔}{\textbf{False.} $4\times 8=32$ and $4\times 9=36$, so $35$ is stepped over.}}

\tf{\ealpara{$3$ کا ہر \ealkeytr{مضاعف} $9$ کا بھی \ealkeytr{مضاعف} ہے}{Every \ealkey{multiple} of $3$ is also a \ealkey{multiple} of $9$}}
  {\ealpara{\textbf{جھوٹ۔} $6$، $3$ کا \ealkeytr{مضاعف} ہے لیکن $9$ کا نہیں۔\\[0.3em] دوسرا طریقہ \emph{سچ} ہے۔}{\textbf{False.} $6$ is a \ealkey{multiple} of $3$ but not of $9$.\\[0.3em] The other way round \emph{is} true.}}

\tf[\Large]{\ealpara{دو اعداد کا \ealkeytr{سب سے چھوٹا مشترک مضاعف} ہمیشہ انہیں آپس میں ضرب دے کر حاصل ہوتا ہے}{The \ealkey{lowest common multiple} of two numbers is always found by multiplying them together}}
  {\ealpara{\textbf{جھوٹ۔} $4\times 6=24$، لیکن $4$ اور $6$ کا \ealkeytr{LCM} $12$ ہے۔\\[0.3em] ضرب دینا صرف اس وقت کام کرتا ہے جب اعداد میں کوئی مشترک \ealkeytr{عامل} نہ ہو۔}{\textbf{False.} $4\times 6=24$, but the \ealkey{LCM} of $4$ and $6$ is $12$.\\[0.3em] Multiplying only works when the numbers share no \ealkey{factors}.}}

\qq[\LARGE]{\ealpara{$4$ کے پہلے پانچ \ealkeytr{مضاعف} لکھیں}{Write down the first five \ealkey{multiples} of $4$}}
  {$4,\ 8,\ 12,\ 16,\ 20$}

\qq[\LARGE]{\ealpara{$7$ کے پہلے پانچ \ealkeytr{مضاعف} لکھیں}{Write down the first five \ealkey{multiples} of $7$}}
  {$7,\ 14,\ 21,\ 28,\ 35$}

\qq[\LARGE]{\ealpara{$20$ اور $50$ کے درمیان $6$ کے \ealkeytr{مضاعف} لکھیں}{Write down the \ealkey{multiples} of $6$ between $20$ and $50$}}
  {$24,\ 30,\ 36,\ 42,\ 48$}

\qq[\LARGE]{\ealpara{کیا $126$، $9$ کا \ealkeytr{مضاعف} ہے؟}{Is $126$ a \ealkey{multiple} of $9$?}}
  {\ealpara{جی ہاں۔ $1+2+6=9$، اور $9\times 14=126$۔}{Yes. $1+2+6=9$, and $9\times 14=126$.}}

\qq[\LARGE]{\ealpara{$4$ اور $6$ کا سب سے چھوٹا مشترک مضاعف تلاش کریں}{Find the \ealkey{lowest common multiple} of $4$ and $6$}}
  {$4,\ 8,\ \mathbf{12},\ 16,\ \ldots$\\[0.3em]
   $6,\ \mathbf{12},\ 18,\ \ldots$\\[0.3em]
   $\mathrm{LCM}=12$}

\qq[\LARGE]{\ealpara{$8$ اور $12$ کا سب سے چھوٹا مشترک مضاعف تلاش کریں}{Find the \ealkey{lowest common multiple} of $8$ and $12$}}
  {$8,\ 16,\ \mathbf{24},\ 32,\ \ldots$\\[0.3em]
   $12,\ \mathbf{24},\ 36,\ \ldots$\\[0.3em]
   $\mathrm{LCM}=24$}

\qq[\LARGE]{\ealpara{$5$، $6$ اور $9$ کا سب سے چھوٹا مشترک مضاعف تلاش کریں}{Find the \ealkey{lowest common multiple} of $5$, $6$ and $9$}}
  {\ealpara{$5$ اور $6$ کا \ealkeytr{سب سے چھوٹا مشترک مضاعف} $30$ ہے۔\\[0.3em]
   $30$ اور $9$ کا \ealkeytr{سب سے چھوٹا مشترک مضاعف} $90$ ہے۔\\[0.3em]
   $\mathrm{LCM}=90$}{\ealkey{LCM} of $5$ and $6$ is $30$.\\[0.3em]
   \ealkey{LCM} of $30$ and $9$ is $90$.\\[0.3em]
   $\mathrm{LCM}=90$}}

\qq[\LARGE]{\ealpara{ایک عدد $4$ اور $10$ دونوں کا \ealkeytr{مضاعف} ہے۔ سب سے چھوٹا کیا ہو سکتا ہے؟}{A number is a \ealkey{multiple} of both $4$ and $10$.\\[0.4em] What is the smallest it could be?}}
  {\ealpara{$20$، $4$ اور $10$ کا سب سے چھوٹا مشترک مضاعف۔}{$20$, the \ealkey{lowest common multiple} of $4$ and $10$.}}

\qq[\Large][\Large]{\ealpara{دو لائٹ ہاؤس اب ایک ساتھ چمکتے ہیں۔ ایک ہر $12$ سیکنڈ میں، دوسرا ہر $18$ سیکنڈ میں۔ وہ اگلی بار ایک ساتھ کب چمکیں گے؟}{Two lighthouses flash together now.\\[0.4em] One flashes every $12$ seconds, the other every $18$ seconds.\\[0.4em] When do they next flash together?}}
  {\ealpara{$\mathrm{LCM}(12,18)=36$\\[0.3em] $36$ سیکنڈ میں۔}{$\mathrm{LCM}(12,18)=36$\\[0.3em] In $36$ seconds.}}

\qq[\Large][\Large]{\ealpara{$1$ سے اوپر سب سے چھوٹا عدد تلاش کریں جو $4$، $5$ اور $6$ سے تقسیم کرنے پر $1$ باقی چھوڑتا ہے}{Find the smallest number above $1$ that leaves a remainder of $1$ when divided by $4$, by $5$ and by $6$}}
  {\ealpara{عدد کو $4$، $5$ اور $6$ کے مشترک \ealkeytr{مضاعف} سے ایک زیادہ ہونا چاہیے۔ $\mathrm{LCM}(4,5,6)=60$، اس لیے عدد $61$ ہے۔}{The number must be one more than a common \ealkey{multiple} of $4$, $5$ and $6$.\\[0.3em] $\mathrm{LCM}(4,5,6)=60$, so the number is $61$.}}

%================================================================
\qtopic{Primes}
%================================================================

\tf{\ealpara{$1$ ایک \ealkeytr{مفرد عدد} ہے}{$1$ is a \ealkey{prime number}}}
  {\ealpara{\textbf{جھوٹ۔} ایک \ealkeytr{مفرد عدد} کے بالکل دو \ealkeytr{عوامل} ہوتے ہیں، لیکن $1$ کے صرف ایک ہے۔}{\textbf{False.} A \ealkey{prime} has exactly two \ealkey{factors}, but $1$ has only one.}}

\tf{\ealpara{$2$ واحد جفت \ealkeytr{مفرد عدد} ہے}{$2$ is the only even \ealkey{prime number}}}
  {\ealpara{\textbf{سچ۔} باقی ہر جفت عدد میں $2$ ایک اضافی \ealkeytr{عامل} کے طور پر ہوتا ہے، اس لیے اس کے دو سے زیادہ \ealkeytr{عوامل} ہوتے ہیں۔}{\textbf{True.} Every other even number has $2$ as an extra \ealkey{factor}, so it has more than two \ealkey{factors}.}}

\tf{\ealpara{$51$ ایک \ealkeytr{مفرد عدد} ہے}{$51$ is a \ealkey{prime number}}}
  {\ealpara{\textbf{جھوٹ۔} $51=3\times 17$۔ اس کے ہندسے $6$ بناتے ہیں، اس لیے یہ $3$ سے تقسیم ہو جاتا ہے۔}{\textbf{False.} $51=3\times 17$.\\[0.3em] Its digits add to $6$, so it divides by $3$.}}

\tf{\ealpara{لامحدود \ealkeytr{مفرد اعداد} ہیں}{There are infinitely many \ealkey{prime numbers}}}
  {\ealpara{\textbf{سچ۔} \ealkeytr{مفرد اعداد} کبھی ختم نہیں ہوتے، حالانکہ اعداد بڑے ہونے پر وہ کم ہوتے جاتے ہیں۔}{\textbf{True.} The \ealkey{primes} never run out, although they do thin out as the numbers get larger.}}

\tf{\ealpara{اگر کوئی عدد $3$ پر ختم ہو، تو وہ لازمی \ealkeytr{مفرد عدد} ہوگا}{If a number ends in $3$, it must be \ealkey{prime}}}
  {\ealpara{\textbf{جھوٹ۔} $33=3\times 11$ اور $63=7\times 9$۔}{\textbf{False.} $33=3\times 11$ and $63=7\times 9$.}}

\qq[\LARGE]{\ealpara{$1$ اور $10$ کے درمیان تمام \ealkeytr{مفرد اعداد} لکھیں}{Write down all the \ealkey{prime numbers} between $1$ and $10$}}
  {$2,\ 3,\ 5,\ 7$}

\qq[\LARGE]{\ealpara{$10$ اور $20$ کے درمیان تمام \ealkeytr{مفرد اعداد} لکھیں}{Write down all the \ealkey{prime numbers} between $10$ and $20$}}
  {$11,\ 13,\ 17,\ 19$}

\qq[\LARGE]{\ealpara{کیا $27$ ایک \ealkeytr{مفرد عدد} ہے؟ اپنا جواب واضح کریں۔}{Is $27$ a \ealkey{prime number}?\\[0.4em]Explain your answer.}}
  {\ealpara{نہیں۔ $27=3\times 9$، اس لیے اس کے دو سے زیادہ \ealkeytr{عوامل} ہیں۔}{No. $27=3\times 9$, so it has more than two \ealkey{factors}.}}

\qq[\LARGE]{\ealpara{$30$ سے بڑا پہلا \ealkeytr{مفرد عدد} لکھیں}{Write down the first \ealkey{prime number} greater than $30$}}
  {$31$}

\qq[\LARGE]{\ealpara{کیا $91$ ایک \ealkeytr{مفرد عدد} ہے؟}{Is $91$ a \ealkey{prime number}?}}
  {\ealpara{نہیں۔ $91=7\times 13$۔}{No. $91=7\times 13$.}}

\qq[\LARGE]{\ealpara{$20$ اور $40$ کے درمیان کتنے \ealkeytr{مفرد اعداد} ہیں؟}{How many \ealkey{prime numbers} are there between $20$ and $40$?}}
  {\ealpara{$23,\ 29,\ 31,\ 37$\\[0.3em] ان میں سے چار۔}{$23,\ 29,\ 31,\ 37$\\[0.3em]Four of them.}}

\qq[\LARGE]{\ealpara{کیا $143$ ایک \ealkeytr{مفرد عدد} ہے؟}{Is $143$ a \ealkey{prime number}?}}
  {\ealpara{نہیں۔ $143=11\times 13$۔}{No. $143=11\times 13$.}}

\qq[\LARGE]{\ealpara{دو \ealkeytr{مفرد اعداد} تلاش کریں جن کا مجموعہ $24$ ہو}{Find two \ealkey{prime numbers} that add up to $24$}}
  {$5+19$ \quad $7+17$ \quad $11+13$}

\qq[\Large][\Large]{\ealpara{یہ جانچنے کے لیے کہ $211$ \ealkeytr{مفرد عدد} ہے، آپ کو کن اعداد سے تقسیم کرنا ہوگا؟ کیا یہ \ealkeytr{مفرد عدد} ہے؟}{To test whether $211$ is \ealkey{prime},\\[0.3em] which numbers do you need to divide by?\\[0.4em]Is it \ealkey{prime}?}}
  {\ealpara{صرف $\sqrt{211}\approx 14.5$ تک کے \ealkeytr{مفرد اعداد}: $2,\ 3,\ 5,\ 7,\ 11,\ 13$۔ ان میں سے کوئی بھی $211$ کو تقسیم نہیں کرتا، اس لیے $211$ \ealkeytr{مفرد عدد} ہے۔}{Only the \ealkey{primes} up to $\sqrt{211}\approx 14.5$:\\[0.3em] $2,\ 3,\ 5,\ 7,\ 11,\ 13$\\[0.3em] None of them divides $211$, so $211$ is \ealkey{prime}.}}

\qq[\Large][\Large]{\ealpara{وضاحت کریں کہ $n^2-1$ جب $n>2$ ہو تو \ealkeytr{مفرد عدد} کیوں نہیں ہو سکتا}{Explain why $n^2-1$ can never be \ealkey{prime} when $n>2$}}
  {\ealpara{$n^2-1=(n-1)(n+1)$\\[0.3em] جب $n>2$ ہو تو دونوں بریکٹ $1$ سے بڑے ہوتے ہیں، اس لیے عدد کے ہمیشہ دو سے زیادہ \ealkeytr{عوامل} ہوتے ہیں۔}{$n^2-1=(n-1)(n+1)$\\[0.3em] When $n>2$ both brackets are bigger than $1$, so the number always has more than two \ealkey{factors}.}}

%================================================================
\qtopic{Is it prime?}
%================================================================

\prq{1}{$1$}
  {\ealpara{$1$ کے صرف ایک \ealkeytr{عامل} ہے۔ ایک \ealkeytr{مفرد عدد} کے بالکل دو مختلف \ealkeytr{عوامل} ہوتے ہیں، $1$ اور خود، اور $1$ کے لیے وہ ایک ہی \ealkeytr{عامل} ہوگا، اس لیے $1$ \ealkeytr{مفرد عدد} نہیں ہے۔}{$1$ has only one \ealkey{factor}. A \ealkey{prime number} has exactly two different \ealkey{factors}, $1$ and itself, and for $1$ those would be the same \ealkey{factor}, so $1$ is not a \ealkey{prime number}.}}

\prq{2}{$1$, $2$}
  {\ealpara{$2$ کے بالکل دو \ealkeytr{عوامل} ہیں، $1$ اور خود، اس لیے $2$ ایک \ealkeytr{مفرد عدد} ہے۔ یہ واحد جفت \ealkeytr{مفرد عدد} ہے، کیونکہ باقی ہر جفت عدد میں $2$ تیسرا \ealkeytr{عامل} ہوتا ہے۔}{$2$ has exactly two \ealkey{factors}, $1$ and itself, so $2$ is a \ealkey{prime number}. It is the only even \ealkey{prime number}, because every other even number has $2$ as a third \ealkey{factor}.}}

\prq{3}{$1$, $3$}
  {\ealpara{$3$ کے بالکل دو \ealkeytr{عوامل} ہیں، $1$ اور خود، اس لیے $3$ ایک \ealkeytr{مفرد عدد} ہے۔}{$3$ has exactly two \ealkey{factors}, $1$ and itself, so $3$ is a \ealkey{prime number}.}}

\prq{4}{$1$, $2$, $4$}
  {\ealpara{$4$ کے تین \ealkeytr{عوامل} ہیں، دو نہیں، اس لیے $4$ \ealkeytr{مفرد عدد} نہیں ہے۔}{$4$ has three \ealkey{factors}, not two, so $4$ is not a \ealkey{prime number}.}}

\prq{5}{$1$, $5$}
  {\ealpara{$5$ کے بالکل دو \ealkeytr{عوامل} ہیں، $1$ اور خود، اس لیے $5$ ایک \ealkeytr{مفرد عدد} ہے۔\\[0.45em] ہمیں صرف $2$ کو \ealkeytr{عامل} کے طور پر جانچنا ہے، کیونکہ $3\times 3=9$، $5$ سے بڑا ہے۔}{$5$ has exactly two \ealkey{factors}, $1$ and itself, so $5$ is a \ealkey{prime number}.\\[0.45em]We only need to test $2$ as a \ealkey{factor}, because $3\times 3=9$ is bigger than $5$.}}

\prq{6}{$1$, $2$, $3$, $6$}
  {\ealpara{$6$ کے چار \ealkeytr{عوامل} ہیں، دو نہیں، اس لیے $6$ \ealkeytr{مفرد عدد} نہیں ہے۔}{$6$ has four \ealkey{factors}, not two, so $6$ is not a \ealkey{prime number}.}}

\prq{7}{$1$, $7$}
  {\ealpara{$7$ کے بالکل دو \ealkeytr{عوامل} ہیں، $1$ اور خود، اس لیے $7$ ایک \ealkeytr{مفرد عدد} ہے۔\\[0.45em] ہمیں صرف $2$ کو \ealkeytr{عامل} کے طور پر جانچنا ہے، کیونکہ $3\times 3=9$، $7$ سے بڑا ہے۔}{$7$ has exactly two \ealkey{factors}, $1$ and itself, so $7$ is a \ealkey{prime number}.\\[0.45em]We only need to test $2$ as a \ealkey{factor}, because $3\times 3=9$ is bigger than $7$.}}

\prq{8}{$1$, $2$, $4$, $8$}
  {\ealpara{$8$ کے چار \ealkeytr{عوامل} ہیں، دو نہیں، اس لیے $8$ \ealkeytr{مفرد عدد} نہیں ہے۔}{$8$ has four \ealkey{factors}, not two, so $8$ is not a \ealkey{prime number}.}}

\prq{9}{$1$, $3$, $9$}
  {\ealpara{$9$ کے تین \ealkeytr{عوامل} ہیں، دو نہیں، اس لیے $9$ \ealkeytr{مفرد عدد} نہیں ہے۔}{$9$ has three \ealkey{factors}, not two, so $9$ is not a \ealkey{prime number}.}}

\prq{10}{$1$, $2$, $5$, $10$}
  {\ealpara{$10$ کے چار \ealkeytr{عوامل} ہیں، دو نہیں، اس لیے $10$ \ealkeytr{مفرد عدد} نہیں ہے۔}{$10$ has four \ealkey{factors}, not two, so $10$ is not a \ealkey{prime number}.}}

\prq{11}{$1$, $11$}
  {\ealpara{$11$ کے بالکل دو \ealkeytr{عوامل} ہیں، $1$ اور خود، اس لیے $11$ ایک \ealkeytr{مفرد عدد} ہے۔\\[0.45em] ہمیں صرف $2$ اور $3$ کو \ealkeytr{عوامل} کے طور پر جانچنا ہے، کیونکہ $5\times 5=25$، $11$ سے بڑا ہے۔}{$11$ has exactly two \ealkey{factors}, $1$ and itself, so $11$ is a \ealkey{prime number}.\\[0.45em]We only need to test $2$ and $3$ as \ealkey{factors}, because $5\times 5=25$ is bigger than $11$.}}

\prq{12}{$1$, $2$, $3$, $4$, $6$, $12$}
  {\ealpara{$12$ کے چھ \ealkeytr{عوامل} ہیں، دو نہیں، اس لیے $12$ \ealkeytr{مفرد عدد} نہیں ہے۔}{$12$ has six \ealkey{factors}, not two, so $12$ is not a \ealkey{prime number}.}}

\prq{13}{$1$, $13$}
  {\ealpara{$13$ کے بالکل دو \ealkeytr{عوامل} ہیں، $1$ اور خود، اس لیے $13$ ایک \ealkeytr{مفرد عدد} ہے۔\\[0.45em] ہمیں صرف $2$ اور $3$ کو \ealkeytr{عوامل} کے طور پر جانچنا ہے، کیونکہ $5\times 5=25$، $13$ سے بڑا ہے۔}{$13$ has exactly two \ealkey{factors}, $1$ and itself, so $13$ is a \ealkey{prime number}.\\[0.45em]We only need to test $2$ and $3$ as \ealkey{factors}, because $5\times 5=25$ is bigger than $13$.}}

\prq{14}{$1$, $2$, $7$, $14$}
  {\ealpara{$14$ کے چار \ealkeytr{عوامل} ہیں، دو نہیں، اس لیے $14$ \ealkeytr{مفرد عدد} نہیں ہے۔}{$14$ has four \ealkey{factors}, not two, so $14$ is not a \ealkey{prime number}.}}

\prq{15}{$1$, $3$, $5$, $15$}
  {\ealpara{$15$ کے چار \ealkeytr{عوامل} ہیں، دو نہیں، اس لیے $15$ \ealkeytr{مفرد عدد} نہیں ہے۔}{$15$ has four \ealkey{factors}, not two, so $15$ is not a \ealkey{prime number}.}}

\prq{16}{$1$, $2$, $4$, $8$, $16$}
  {\ealpara{$16$ کے پانچ \ealkeytr{عوامل} ہیں، دو نہیں، اس لیے $16$ \ealkeytr{مفرد عدد} نہیں ہے۔}{$16$ has five \ealkey{factors}, not two, so $16$ is not a \ealkey{prime number}.}}

\prq{17}{$1$, $17$}
  {\ealpara{$17$ کے بالکل دو \ealkeytr{عوامل} ہیں، $1$ اور خود، اس لیے $17$ ایک \ealkeytr{مفرد عدد} ہے۔\\[0.45em] ہمیں صرف $2$ اور $3$ کو \ealkeytr{عوامل} کے طور پر جانچنا ہے، کیونکہ $5\times 5=25$، $17$ سے بڑا ہے۔}{$17$ has exactly two \ealkey{factors}, $1$ and itself, so $17$ is a \ealkey{prime number}.\\[0.45em]We only need to test $2$ and $3$ as \ealkey{factors}, because $5\times 5=25$ is bigger than $17$.}}

\prq{18}{$1$, $2$, $3$, $6$, $9$, $18$}
  {\ealpara{$18$ کے چھ \ealkeytr{عوامل} ہیں، دو نہیں، اس لیے $18$ \ealkeytr{مفرد عدد} نہیں ہے۔}{$18$ has six \ealkey{factors}, not two, so $18$ is not a \ealkey{prime number}.}}

\prq{19}{$1$, $19$}
  {\ealpara{$19$ کے بالکل دو \ealkeytr{عوامل} ہیں، $1$ اور خود، اس لیے $19$ ایک \ealkeytr{مفرد عدد} ہے۔\\[0.45em] ہمیں صرف $2$ اور $3$ کو \ealkeytr{عوامل} کے طور پر جانچنا ہے، کیونکہ $5\times 5=25$، $19$ سے بڑا ہے۔}{$19$ has exactly two \ealkey{factors}, $1$ and itself, so $19$ is a \ealkey{prime number}.\\[0.45em]We only need to test $2$ and $3$ as \ealkey{factors}, because $5\times 5=25$ is bigger than $19$.}}

\prq{20}{$1$, $2$, $4$, $5$, $10$, $20$}
  {\ealpara{$20$ کے چھ \ealkeytr{عوامل} ہیں، دو نہیں، اس لیے $20$ \ealkeytr{مفرد عدد} نہیں ہے۔}{$20$ has six \ealkey{factors}, not two, so $20$ is not a \ealkey{prime number}.}}

\prq{21}{$1$, $3$, $7$, $21$}
  {\ealpara{$21$ کے چار \ealkeytr{عوامل} ہیں، دو نہیں، اس لیے $21$ \ealkeytr{مفرد عدد} نہیں ہے۔}{$21$ has four \ealkey{factors}, not two, so $21$ is not a \ealkey{prime number}.}}

\prq{22}{$1$, $2$, $11$, $22$}
  {\ealpara{$22$ کے چار \ealkeytr{عوامل} ہیں، دو نہیں، اس لیے $22$ \ealkeytr{مفرد عدد} نہیں ہے۔}{$22$ has four \ealkey{factors}, not two, so $22$ is not a \ealkey{prime number}.}}

\prq{23}{$1$, $23$}
  {\ealpara{$23$ کے بالکل دو \ealkeytr{عوامل} ہیں، $1$ اور خود، اس لیے $23$ ایک \ealkeytr{مفرد عدد} ہے۔\\[0.45em] ہمیں صرف $2$ اور $3$ کو \ealkeytr{عوامل} کے طور پر جانچنا ہے، کیونکہ $5\times 5=25$، $23$ سے بڑا ہے۔}{$23$ has exactly two \ealkey{factors}, $1$ and itself, so $23$ is a \ealkey{prime number}.\\[0.45em]We only need to test $2$ and $3$ as \ealkey{factors}, because $5\times 5=25$ is bigger than $23$.}}

\prq{24}{$1$, $2$, $3$, $4$, $6$, $8$, $12$, $24$}
  {\ealpara{$24$ کے آٹھ \ealkeytr{عوامل} ہیں، دو نہیں، اس لیے $24$ \ealkeytr{مفرد عدد} نہیں ہے۔}{$24$ has eight \ealkey{factors}, not two, so $24$ is not a \ealkey{prime number}.}}

\prq{25}{$1$, $5$, $25$}
  {\ealpara{$25$ کے تین \ealkeytr{عوامل} ہیں، دو نہیں، اس لیے $25$ \ealkeytr{مفرد عدد} نہیں ہے۔}{$25$ has three \ealkey{factors}, not two, so $25$ is not a \ealkey{prime number}.}}

\prq{26}{$1$, $2$, $13$, $26$}
  {\ealpara{$26$ کے چار \ealkeytr{عوامل} ہیں، دو نہیں، اس لیے $26$ \ealkeytr{مفرد عدد} نہیں ہے۔}{$26$ has four \ealkey{factors}, not two, so $26$ is not a \ealkey{prime number}.}}

\prq{27}{$1$, $3$, $9$, $27$}
  {\ealpara{$27$ کے چار \ealkeytr{عوامل} ہیں، دو نہیں، اس لیے $27$ \ealkeytr{مفرد عدد} نہیں ہے۔}{$27$ has four \ealkey{factors}, not two, so $27$ is not a \ealkey{prime number}.}}

\prq{28}{$1$, $2$, $4$, $7$, $14$, $28$}
  {\ealpara{$28$ کے چھ \ealkeytr{عوامل} ہیں، دو نہیں، اس لیے $28$ \ealkeytr{مفرد عدد} نہیں ہے۔}{$28$ has six \ealkey{factors}, not two, so $28$ is not a \ealkey{prime number}.}}

\prq{29}{$1$, $29$}
  {\ealpara{$29$ کے بالکل دو \ealkeytr{عوامل} ہیں، $1$ اور خود، اس لیے $29$ ایک \ealkeytr{مفرد عدد} ہے۔\\[0.45em] ہمیں صرف $2$، $3$ اور $5$ کو \ealkeytr{عوامل} کے طور پر جانچنا ہے، کیونکہ $7\times 7=49$، $29$ سے بڑا ہے۔}{$29$ has exactly two \ealkey{factors}, $1$ and itself, so $29$ is a \ealkey{prime number}.\\[0.45em]We only need to test $2$, $3$ and $5$ as \ealkey{factors}, because $7\times 7=49$ is bigger than $29$.}}

\prq{30}{$1$, $2$, $3$, $5$, $6$, $10$, $15$, $30$}
  {\ealpara{$30$ کے آٹھ \ealkeytr{عوامل} ہیں، دو نہیں، اس لیے $30$ \ealkeytr{مفرد عدد} نہیں ہے۔}{$30$ has eight \ealkey{factors}, not two, so $30$ is not a \ealkey{prime number}.}}

\prq{31}{$1$, $31$}
  {\ealpara{$31$ کے بالکل دو \ealkeytr{عوامل} ہیں، $1$ اور خود، اس لیے $31$ ایک \ealkeytr{مفرد عدد} ہے۔\\[0.45em] ہمیں صرف $2$، $3$ اور $5$ کو \ealkeytr{عوامل} کے طور پر جانچنا ہے، کیونکہ $7\times 7=49$، $31$ سے بڑا ہے۔}{$31$ has exactly two \ealkey{factors}, $1$ and itself, so $31$ is a \ealkey{prime number}.\\[0.45em]We only need to test $2$, $3$ and $5$ as \ealkey{factors}, because $7\times 7=49$ is bigger than $31$.}}

\prq{32}{$1$, $2$, $4$, $8$, $16$, $32$}
  {\ealpara{$32$ کے چھ \ealkeytr{عوامل} ہیں، دو نہیں، اس لیے $32$ \ealkeytr{مفرد عدد} نہیں ہے۔}{$32$ has six \ealkey{factors}, not two, so $32$ is not a \ealkey{prime number}.}}

\prq{33}{$1$, $3$, $11$, $33$}
  {\ealpara{$33$ کے چار \ealkeytr{عوامل} ہیں، دو نہیں، اس لیے $33$ \ealkeytr{مفرد عدد} نہیں ہے۔}{$33$ has four \ealkey{factors}, not two, so $33$ is not a \ealkey{prime number}.}}

\prq{34}{$1$, $2$, $17$, $34$}
  {\ealpara{$34$ کے چار \ealkeytr{عوامل} ہیں، دو نہیں، اس لیے $34$ \ealkeytr{مفرد عدد} نہیں ہے۔}{$34$ has four \ealkey{factors}, not two, so $34$ is not a \ealkey{prime number}.}}

\prq{35}{$1$, $5$, $7$, $35$}
  {\ealpara{$35$ کے چار \ealkeytr{عوامل} ہیں، دو نہیں، اس لیے $35$ \ealkeytr{مفرد عدد} نہیں ہے۔}{$35$ has four \ealkey{factors}, not two, so $35$ is not a \ealkey{prime number}.}}

\prq{36}{$1$, $2$, $3$, $4$, $6$, $9$, $12$, $18$, $36$}
  {\ealpara{$36$ کے نو \ealkeytr{عوامل} ہیں، دو نہیں، اس لیے $36$ \ealkeytr{مفرد عدد} نہیں ہے۔}{$36$ has nine \ealkey{factors}, not two, so $36$ is not a \ealkey{prime number}.}}

\prq{37}{$1$, $37$}
  {\ealpara{$37$ کے بالکل دو \ealkeytr{عوامل} ہیں، $1$ اور خود، اس لیے $37$ ایک \ealkeytr{مفرد عدد} ہے۔\\[0.45em] ہمیں صرف $2$، $3$ اور $5$ کو \ealkeytr{عوامل} کے طور پر جانچنا ہے، کیونکہ $7\times 7=49$، $37$ سے بڑا ہے۔}{$37$ has exactly two \ealkey{factors}, $1$ and itself, so $37$ is a \ealkey{prime number}.\\[0.45em]We only need to test $2$, $3$ and $5$ as \ealkey{factors}, because $7\times 7=49$ is bigger than $37$.}}

\prq{38}{$1$, $2$, $19$, $38$}
  {\ealpara{$38$ کے چار \ealkeytr{عوامل} ہیں، دو نہیں، اس لیے $38$ \ealkeytr{مفرد عدد} نہیں ہے۔}{$38$ has four \ealkey{factors}, not two, so $38$ is not a \ealkey{prime number}.}}

\prq{39}{$1$, $3$, $13$, $39$}
  {\ealpara{$39$ کے چار \ealkeytr{عوامل} ہیں، دو نہیں، اس لیے $39$ \ealkeytr{مفرد عدد} نہیں ہے۔}{$39$ has four \ealkey{factors}, not two, so $39$ is not a \ealkey{prime number}.}}

\prq{40}{$1$, $2$, $4$, $5$, $8$, $10$, $20$, $40$}
  {\ealpara{$40$ کے آٹھ \ealkeytr{عوامل} ہیں، دو نہیں، اس لیے $40$ \ealkeytr{مفرد عدد} نہیں ہے۔}{$40$ has eight \ealkey{factors}, not two, so $40$ is not a \ealkey{prime number}.}}

\prq{41}{$1$, $41$}
  {\ealpara{$41$ کے بالکل دو \ealkeytr{عوامل} ہیں، $1$ اور خود، اس لیے $41$ ایک \ealkeytr{مفرد عدد} ہے۔\\[0.45em] ہمیں صرف $2$، $3$ اور $5$ کو \ealkeytr{عوامل} کے طور پر جانچنا ہے، کیونکہ $7\times 7=49$، $41$ سے بڑا ہے۔}{$41$ has exactly two \ealkey{factors}, $1$ and itself, so $41$ is a \ealkey{prime number}.\\[0.45em]We only need to test $2$, $3$ and $5$ as \ealkey{factors}, because $7\times 7=49$ is bigger than $41$.}}

\prq{42}{$1$, $2$, $3$, $6$, $7$, $14$, $21$, $42$}
  {\ealpara{$42$ کے آٹھ \ealkeytr{عوامل} ہیں، دو نہیں، اس لیے $42$ \ealkeytr{مفرد عدد} نہیں ہے۔}{$42$ has eight \ealkey{factors}, not two, so $42$ is not a \ealkey{prime number}.}}

\prq{43}{$1$, $43$}
  {\ealpara{$43$ کے بالکل دو \ealkeytr{عوامل} ہیں، $1$ اور خود، اس لیے $43$ ایک \ealkeytr{مفرد عدد} ہے۔\\[0.45em] ہمیں صرف $2$، $3$ اور $5$ کو \ealkeytr{عوامل} کے طور پر جانچنا ہے، کیونکہ $7\times 7=49$، $43$ سے بڑا ہے۔}{$43$ has exactly two \ealkey{factors}, $1$ and itself, so $43$ is a \ealkey{prime number}.\\[0.45em]We only need to test $2$, $3$ and $5$ as \ealkey{factors}, because $7\times 7=49$ is bigger than $43$.}}

\prq{44}{$1$, $2$, $4$, $11$, $22$, $44$}
  {\ealpara{$44$ کے چھ \ealkeytr{عوامل} ہیں، دو نہیں، اس لیے $44$ \ealkeytr{مفرد عدد} نہیں ہے۔}{$44$ has six \ealkey{factors}, not two, so $44$ is not a \ealkey{prime number}.}}

\prq{45}{$1$, $3$, $5$, $9$, $15$, $45$}
  {\ealpara{$45$ کے چھ \ealkeytr{عوامل} ہیں، دو نہیں، اس لیے $45$ \ealkeytr{مفرد عدد} نہیں ہے۔}{$45$ has six \ealkey{factors}, not two, so $45$ is not a \ealkey{prime number}.}}

\prq{46}{$1$, $2$, $23$, $46$}
  {\ealpara{$46$ کے چار \ealkeytr{عوامل} ہیں، دو نہیں، اس لیے $46$ \ealkeytr{مفرد عدد} نہیں ہے۔}{$46$ has four \ealkey{factors}, not two, so $46$ is not a \ealkey{prime number}.}}

\prq{47}{$1$, $47$}
  {\ealpara{$47$ کے بالکل دو \ealkeytr{عوامل} ہیں، $1$ اور خود، اس لیے $47$ ایک \ealkeytr{مفرد عدد} ہے۔\\[0.45em] ہمیں صرف $2$، $3$ اور $5$ کو \ealkeytr{عوامل} کے طور پر جانچنا ہے، کیونکہ $7\times 7=49$، $47$ سے بڑا ہے۔}{$47$ has exactly two \ealkey{factors}, $1$ and itself, so $47$ is a \ealkey{prime number}.\\[0.45em]We only need to test $2$, $3$ and $5$ as \ealkey{factors}, because $7\times 7=49$ is bigger than $47$.}}

\prq{48}{$1$, $2$, $3$, $4$, $6$, $8$, $12$, $16$, $24$, $48$}
  {\ealpara{$48$ کے دس \ealkeytr{عوامل} ہیں، دو نہیں، اس لیے $48$ \ealkeytr{مفرد عدد} نہیں ہے۔}{$48$ has ten \ealkey{factors}, not two, so $48$ is not a \ealkey{prime number}.}}

\prq{49}{$1$, $7$, $49$}
  {\ealpara{$49$ کے تین \ealkeytr{عوامل} ہیں، دو نہیں، اس لیے $49$ \ealkeytr{مفرد عدد} نہیں ہے۔}{$49$ has three \ealkey{factors}, not two, so $49$ is not a \ealkey{prime number}.}}

\prq{50}{$1$, $2$, $5$, $10$, $25$, $50$}
  {\ealpara{$50$ کے چھ \ealkeytr{عوامل} ہیں، دو نہیں، اس لیے $50$ \ealkeytr{مفرد عدد} نہیں ہے۔}{$50$ has six \ealkey{factors}, not two, so $50$ is not a \ealkey{prime number}.}}

\prq{51}{$1$, $3$, $17$, $51$}
  {\ealpara{$51$ کے چار \ealkeytr{عوامل} ہیں، دو نہیں، اس لیے $51$ \ealkeytr{مفرد عدد} نہیں ہے۔}{$51$ has four \ealkey{factors}, not two, so $51$ is not a \ealkey{prime number}.}}

\prq{52}{$1$, $2$, $4$, $13$, $26$, $52$}
  {\ealpara{$52$ کے چھ \ealkeytr{عوامل} ہیں، دو نہیں، اس لیے $52$ \ealkeytr{مفرد عدد} نہیں ہے۔}{$52$ has six \ealkey{factors}, not two, so $52$ is not a \ealkey{prime number}.}}

\prq{53}{$1$, $53$}
  {\ealpara{$53$ کے بالکل دو \ealkeytr{عوامل} ہیں، $1$ اور خود، اس لیے $53$ ایک \ealkeytr{مفرد عدد} ہے۔\\[0.45em] ہمیں صرف $2$، $3$، $5$ اور $7$ کو \ealkeytr{عوامل} کے طور پر جانچنا ہے، کیونکہ $11\times 11=121$، $53$ سے بڑا ہے۔}{$53$ has exactly two \ealkey{factors}, $1$ and itself, so $53$ is a \ealkey{prime number}.\\[0.45em]We only need to test $2$, $3$, $5$ and $7$ as \ealkey{factors}, because $11\times 11=121$ is bigger than $53$.}}

\prq{54}{$1$, $2$, $3$, $6$, $9$, $18$, $27$, $54$}
  {\ealpara{$54$ کے آٹھ \ealkeytr{عوامل} ہیں، دو نہیں، اس لیے $54$ \ealkeytr{مفرد عدد} نہیں ہے۔}{$54$ has eight \ealkey{factors}, not two, so $54$ is not a \ealkey{prime number}.}}

\prq{55}{$1$, $5$, $11$, $55$}
  {\ealpara{$55$ کے چار \ealkeytr{عوامل} ہیں، دو نہیں، اس لیے $55$ \ealkeytr{مفرد عدد} نہیں ہے۔}{$55$ has four \ealkey{factors}, not two, so $55$ is not a \ealkey{prime number}.}}

\prq{56}{$1$, $2$, $4$, $7$, $8$, $14$, $28$, $56$}
  {\ealpara{$56$ کے آٹھ \ealkeytr{عوامل} ہیں، دو نہیں، اس لیے $56$ \ealkeytr{مفرد عدد} نہیں ہے۔}{$56$ has eight \ealkey{factors}, not two, so $56$ is not a \ealkey{prime number}.}}

\prq{57}{$1$, $3$, $19$, $57$}
  {\ealpara{$57$ کے چار \ealkeytr{عوامل} ہیں، دو نہیں، اس لیے $57$ \ealkeytr{مفرد عدد} نہیں ہے۔}{$57$ has four \ealkey{factors}, not two, so $57$ is not a \ealkey{prime number}.}}

\prq{58}{$1$, $2$, $29$, $58$}
  {\ealpara{$58$ کے چار \ealkeytr{عوامل} ہیں، دو نہیں، اس لیے $58$ \ealkeytr{مفرد عدد} نہیں ہے۔}{$58$ has four \ealkey{factors}, not two, so $58$ is not a \ealkey{prime number}.}}

\prq{59}{$1$, $59$}
  {\ealpara{$59$ کے بالکل دو \ealkeytr{عوامل} ہیں، $1$ اور خود، اس لیے $59$ ایک \ealkeytr{مفرد عدد} ہے۔\\[0.45em] ہمیں صرف $2$، $3$، $5$ اور $7$ کو \ealkeytr{عوامل} کے طور پر جانچنا ہے، کیونکہ $11\times 11=121$، $59$ سے بڑا ہے۔}{$59$ has exactly two \ealkey{factors}, $1$ and itself, so $59$ is a \ealkey{prime number}.\\[0.45em]We only need to test $2$, $3$, $5$ and $7$ as \ealkey{factors}, because $11\times 11=121$ is bigger than $59$.}}

\prq{60}{$1$, $2$, $3$, $4$, $5$, $6$, $10$, $12$, $15$, $20$, $30$, $60$}
  {\ealpara{$60$ کے بارہ \ealkeytr{عوامل} ہیں، دو نہیں، اس لیے $60$ \ealkeytr{مفرد عدد} نہیں ہے۔}{$60$ has twelve \ealkey{factors}, not two, so $60$ is not a \ealkey{prime number}.}}

\prq{61}{$1$, $61$}
  {\ealpara{$61$ کے بالکل دو \ealkeytr{عوامل} ہیں، $1$ اور خود، اس لیے $61$ ایک \ealkeytr{مفرد عدد} ہے۔\\[0.45em] ہمیں صرف $2$، $3$، $5$ اور $7$ کو \ealkeytr{عوامل} کے طور پر جانچنا ہے، کیونکہ $11\times 11=121$، $61$ سے بڑا ہے۔}{$61$ has exactly two \ealkey{factors}, $1$ and itself, so $61$ is a \ealkey{prime number}.\\[0.45em]We only need to test $2$, $3$, $5$ and $7$ as \ealkey{factors}, because $11\times 11=121$ is bigger than $61$.}}

\prq{62}{$1$, $2$, $31$, $62$}
  {\ealpara{$62$ کے چار \ealkeytr{عوامل} ہیں، دو نہیں، اس لیے $62$ \ealkeytr{مفرد عدد} نہیں ہے۔}{$62$ has four \ealkey{factors}, not two, so $62$ is not a \ealkey{prime number}.}}

\prq{63}{$1$, $3$, $7$, $9$, $21$, $63$}
  {\ealpara{$63$ کے چھ \ealkeytr{عوامل} ہیں، دو نہیں، اس لیے $63$ \ealkeytr{مفرد عدد} نہیں ہے۔}{$63$ has six \ealkey{factors}, not two, so $63$ is not a \ealkey{prime number}.}}

\prq{64}{$1$, $2$, $4$, $8$, $16$, $32$, $64$}
  {\ealpara{$64$ کے سات \ealkeytr{عوامل} ہیں، دو نہیں، اس لیے $64$ \ealkeytr{مفرد عدد} نہیں ہے۔}{$64$ has seven \ealkey{factors}, not two, so $64$ is not a \ealkey{prime number}.}}

\prq{65}{$1$, $5$, $13$, $65$}
  {\ealpara{$65$ کے چار \ealkeytr{عوامل} ہیں، دو نہیں، اس لیے $65$ \ealkeytr{مفرد عدد} نہیں ہے۔}{$65$ has four \ealkey{factors}, not two, so $65$ is not a \ealkey{prime number}.}}

\prq{66}{$1$, $2$, $3$, $6$, $11$, $22$, $33$, $66$}
  {\ealpara{$66$ کے آٹھ \ealkeytr{عوامل} ہیں، دو نہیں، اس لیے $66$ \ealkeytr{مفرد عدد} نہیں ہے۔}{$66$ has eight \ealkey{factors}, not two, so $66$ is not a \ealkey{prime number}.}}

\prq{67}{$1$, $67$}
  {\ealpara{$67$ کے بالکل دو \ealkeytr{عوامل} ہیں، $1$ اور خود، اس لیے $67$ ایک \ealkeytr{مفرد عدد} ہے۔\\[0.45em] ہمیں صرف $2$، $3$، $5$ اور $7$ کو \ealkeytr{عوامل} کے طور پر جانچنا ہے، کیونکہ $11\times 11=121$، $67$ سے بڑا ہے۔}{$67$ has exactly two \ealkey{factors}, $1$ and itself, so $67$ is a \ealkey{prime number}.\\[0.45em]We only need to test $2$, $3$, $5$ and $7$ as \ealkey{factors}, because $11\times 11=121$ is bigger than $67$.}}

\prq{68}{$1$, $2$, $4$, $17$, $34$, $68$}
  {\ealpara{$68$ کے چھ \ealkeytr{عوامل} ہیں، دو نہیں، اس لیے $68$ \ealkeytr{مفرد عدد} نہیں ہے۔}{$68$ has six \ealkey{factors}, not two, so $68$ is not a \ealkey{prime number}.}}

\prq{69}{$1$, $3$, $23$, $69$}
  {\ealpara{$69$ کے چار \ealkeytr{عوامل} ہیں، دو نہیں، اس لیے $69$ \ealkeytr{مفرد عدد} نہیں ہے۔}{$69$ has four \ealkey{factors}, not two, so $69$ is not a \ealkey{prime number}.}}

\prq{70}{$1$, $2$, $5$, $7$, $10$, $14$, $35$, $70$}
  {\ealpara{$70$ کے آٹھ \ealkeytr{عوامل} ہیں، دو نہیں، اس لیے $70$ \ealkeytr{مفرد عدد} نہیں ہے۔}{$70$ has eight \ealkey{factors}, not two, so $70$ is not a \ealkey{prime number}.}}

\prq{71}{$1$, $71$}
  {\ealpara{$71$ کے بالکل دو \ealkeytr{عوامل} ہیں، $1$ اور خود، اس لیے $71$ ایک \ealkeytr{مفرد عدد} ہے۔\\[0.45em] ہمیں صرف $2$، $3$، $5$ اور $7$ کو \ealkeytr{عوامل} کے طور پر جانچنا ہے، کیونکہ $11\times 11=121$، $71$ سے بڑا ہے۔}{$71$ has exactly two \ealkey{factors}, $1$ and itself, so $71$ is a \ealkey{prime number}.\\[0.45em]We only need to test $2$, $3$, $5$ and $7$ as \ealkey{factors}, because $11\times 11=121$ is bigger than $71$.}}

\prq{72}{$1$, $2$, $3$, $4$, $6$, $8$, $9$, $12$, $18$, $24$, $36$, $72$}
  {\ealpara{$72$ کے بارہ \ealkeytr{عوامل} ہیں، دو نہیں، اس لیے $72$ \ealkeytr{مفرد عدد} نہیں ہے۔}{$72$ has twelve \ealkey{factors}, not two, so $72$ is not a \ealkey{prime number}.}}

\prq{73}{$1$, $73$}
  {\ealpara{$73$ کے بالکل دو \ealkeytr{عوامل} ہیں، $1$ اور خود، اس لیے $73$ ایک \ealkeytr{مفرد عدد} ہے۔\\[0.45em] ہمیں صرف $2$، $3$، $5$ اور $7$ کو \ealkeytr{عوامل} کے طور پر جانچنا ہے، کیونکہ $11\times 11=121$، $73$ سے بڑا ہے۔}{$73$ has exactly two \ealkey{factors}, $1$ and itself, so $73$ is a \ealkey{prime number}.\\[0.45em]We only need to test $2$, $3$, $5$ and $7$ as \ealkey{factors}, because $11\times 11=121$ is bigger than $73$.}}

\prq{74}{$1$, $2$, $37$, $74$}
  {\ealpara{$74$ کے چار \ealkeytr{عوامل} ہیں، دو نہیں، اس لیے $74$ \ealkeytr{مفرد عدد} نہیں ہے۔}{$74$ has four \ealkey{factors}, not two, so $74$ is not a \ealkey{prime number}.}}

\prq{75}{$1$, $3$, $5$, $15$, $25$, $75$}
  {\ealpara{$75$ کے چھ \ealkeytr{عوامل} ہیں، دو نہیں، اس لیے $75$ \ealkeytr{مفرد عدد} نہیں ہے۔}{$75$ has six \ealkey{factors}, not two, so $75$ is not a \ealkey{prime number}.}}

\prq{76}{$1$, $2$, $4$, $19$, $38$, $76$}
  {\ealpara{$76$ کے چھ \ealkeytr{عوامل} ہیں، دو نہیں، اس لیے $76$ \ealkeytr{مفرد عدد} نہیں ہے۔}{$76$ has six \ealkey{factors}, not two, so $76$ is not a \ealkey{prime number}.}}

\prq{77}{$1$, $7$, $11$, $77$}
  {\ealpara{$77$ کے چار \ealkeytr{عوامل} ہیں، دو نہیں، اس لیے $77$ \ealkeytr{مفرد عدد} نہیں ہے۔}{$77$ has four \ealkey{factors}, not two, so $77$ is not a \ealkey{prime number}.}}

\prq{78}{$1$, $2$, $3$, $6$, $13$, $26$, $39$, $78$}
  {\ealpara{$78$ کے آٹھ \ealkeytr{عوامل} ہیں، دو نہیں، اس لیے $78$ \ealkeytr{مفرد عدد} نہیں ہے۔}{$78$ has eight \ealkey{factors}, not two, so $78$ is not a \ealkey{prime number}.}}

\prq{79}{$1$, $79$}
  {\ealpara{$79$ کے بالکل دو \ealkeytr{عوامل} ہیں، $1$ اور خود، اس لیے $79$ ایک \ealkeytr{مفرد عدد} ہے۔\\[0.45em] ہمیں صرف $2$، $3$، $5$ اور $7$ کو \ealkeytr{عوامل} کے طور پر جانچنا ہے، کیونکہ $11\times 11=121$، $79$ سے بڑا ہے۔}{$79$ has exactly two \ealkey{factors}, $1$ and itself, so $79$ is a \ealkey{prime number}.\\[0.45em]We only need to test $2$, $3$, $5$ and $7$ as \ealkey{factors}, because $11\times 11=121$ is bigger than $79$.}}

\prq{80}{$1$, $2$, $4$, $5$, $8$, $10$, $16$, $20$, $40$, $80$}
  {\ealpara{$80$ کے دس \ealkeytr{عوامل} ہیں، دو نہیں، اس لیے $80$ \ealkeytr{مفرد عدد} نہیں ہے۔}{$80$ has ten \ealkey{factors}, not two, so $80$ is not a \ealkey{prime number}.}}

\prq{81}{$1$, $3$, $9$, $27$, $81$}
  {\ealpara{$81$ کے پانچ \ealkeytr{عوامل} ہیں، دو نہیں، اس لیے $81$ \ealkeytr{مفرد عدد} نہیں ہے۔}{$81$ has five \ealkey{factors}, not two, so $81$ is not a \ealkey{prime number}.}}

\prq{82}{$1$, $2$, $41$, $82$}
  {\ealpara{$82$ کے چار \ealkeytr{عوامل} ہیں، دو نہیں، اس لیے $82$ \ealkeytr{مفرد عدد} نہیں ہے۔}{$82$ has four \ealkey{factors}, not two, so $82$ is not a \ealkey{prime number}.}}

\prq{83}{$1$, $83$}
  {\ealpara{$83$ کے بالکل دو \ealkeytr{عوامل} ہیں، $1$ اور خود، اس لیے $83$ ایک \ealkeytr{مفرد عدد} ہے۔\\[0.45em] ہمیں صرف $2$، $3$، $5$ اور $7$ کو \ealkeytr{عوامل} کے طور پر جانچنا ہے، کیونکہ $11\times 11=121$، $83$ سے بڑا ہے۔}{$83$ has exactly two \ealkey{factors}, $1$ and itself, so $83$ is a \ealkey{prime number}.\\[0.45em]We only need to test $2$, $3$, $5$ and $7$ as \ealkey{factors}, because $11\times 11=121$ is bigger than $83$.}}

\prq{84}{$1$, $2$, $3$, $4$, $6$, $7$, $12$, $14$, $21$, $28$, $42$, $84$}
  {\ealpara{$84$ کے بارہ \ealkeytr{عوامل} ہیں، دو نہیں، اس لیے $84$ \ealkeytr{مفرد عدد} نہیں ہے۔}{$84$ has twelve \ealkey{factors}, not two, so $84$ is not a \ealkey{prime number}.}}

\prq{85}{$1$, $5$, $17$, $85$}
  {\ealpara{$85$ کے چار \ealkeytr{عوامل} ہیں، دو نہیں، اس لیے $85$ \ealkeytr{مفرد عدد} نہیں ہے۔}{$85$ has four \ealkey{factors}, not two, so $85$ is not a \ealkey{prime number}.}}

\prq{86}{$1$, $2$, $43$, $86$}
  {\ealpara{$86$ کے چار \ealkeytr{عوامل} ہیں، دو نہیں، اس لیے $86$ \ealkeytr{مفرد عدد} نہیں ہے۔}{$86$ has four \ealkey{factors}, not two, so $86$ is not a \ealkey{prime number}.}}

\prq{87}{$1$, $3$, $29$, $87$}
  {\ealpara{$87$ کے چار \ealkeytr{عوامل} ہیں، دو نہیں، اس لیے $87$ \ealkeytr{مفرد عدد} نہیں ہے۔}{$87$ has four \ealkey{factors}, not two, so $87$ is not a \ealkey{prime number}.}}

\prq{88}{$1$, $2$, $4$, $8$, $11$, $22$, $44$, $88$}
  {\ealpara{$88$ کے آٹھ \ealkeytr{عوامل} ہیں، دو نہیں، اس لیے $88$ \ealkeytr{مفرد عدد} نہیں ہے۔}{$88$ has eight \ealkey{factors}, not two, so $88$ is not a \ealkey{prime number}.}}

\prq{89}{$1$, $89$}
  {\ealpara{$89$ کے بالکل دو \ealkeytr{عوامل} ہیں، $1$ اور خود، اس لیے $89$ ایک \ealkeytr{مفرد عدد} ہے۔\\[0.45em] ہمیں صرف $2$، $3$، $5$ اور $7$ کو \ealkeytr{عوامل} کے طور پر جانچنا ہے، کیونکہ $11\times 11=121$، $89$ سے بڑا ہے۔}{$89$ has exactly two \ealkey{factors}, $1$ and itself, so $89$ is a \ealkey{prime number}.\\[0.45em]We only need to test $2$, $3$, $5$ and $7$ as \ealkey{factors}, because $11\times 11=121$ is bigger than $89$.}}

\prq{90}{$1$, $2$, $3$, $5$, $6$, $9$, $10$, $15$, $18$, $30$, $45$, $90$}
  {\ealpara{$90$ کے بارہ \ealkeytr{عوامل} ہیں، دو نہیں، اس لیے $90$ \ealkeytr{مفرد عدد} نہیں ہے۔}{$90$ has twelve \ealkey{factors}, not two, so $90$ is not a \ealkey{prime number}.}}

\prq{91}{$1$, $7$, $13$, $91$}
  {\ealpara{$91$ کے چار \ealkeytr{عوامل} ہیں، دو نہیں، اس لیے $91$ \ealkeytr{مفرد عدد} نہیں ہے۔}{$91$ has four \ealkey{factors}, not two, so $91$ is not a \ealkey{prime number}.}}

\prq{92}{$1$, $2$, $4$, $23$, $46$, $92$}
  {\ealpara{$92$ کے چھ \ealkeytr{عوامل} ہیں، دو نہیں، اس لیے $92$ \ealkeytr{مفرد عدد} نہیں ہے۔}{$92$ has six \ealkey{factors}, not two, so $92$ is not a \ealkey{prime number}.}}

\prq{93}{$1$, $3$, $31$, $93$}
  {\ealpara{$93$ کے چار \ealkeytr{عوامل} ہیں، دو نہیں، اس لیے $93$ \ealkeytr{مفرد عدد} نہیں ہے۔}{$93$ has four \ealkey{factors}, not two, so $93$ is not a \ealkey{prime number}.}}

\prq{94}{$1$, $2$, $47$, $94$}
  {\ealpara{$94$ کے چار \ealkeytr{عوامل} ہیں، دو نہیں، اس لیے $94$ \ealkeytr{مفرد عدد} نہیں ہے۔}{$94$ has four \ealkey{factors}, not two, so $94$ is not a \ealkey{prime number}.}}

\prq{95}{$1$, $5$, $19$, $95$}
  {\ealpara{$95$ کے چار \ealkeytr{عوامل} ہیں، دو نہیں، اس لیے $95$ \ealkeytr{مفرد عدد} نہیں ہے۔}{$95$ has four \ealkey{factors}, not two, so $95$ is not a \ealkey{prime number}.}}

\prq{96}{$1$, $2$, $3$, $4$, $6$, $8$, $12$, $16$, $24$, $32$, $48$, $96$}
  {\ealpara{$96$ کے بارہ \ealkeytr{عوامل} ہیں، دو نہیں، اس لیے $96$ \ealkeytr{مفرد عدد} نہیں ہے۔}{$96$ has twelve \ealkey{factors}, not two, so $96$ is not a \ealkey{prime number}.}}

\prq{97}{$1$, $97$}
  {\ealpara{$97$ کے بالکل دو \ealkeytr{عوامل} ہیں، $1$ اور خود، اس لیے $97$ ایک \ealkeytr{مفرد عدد} ہے۔\\[0.45em] ہمیں صرف $2$، $3$، $5$ اور $7$ کو \ealkeytr{عوامل} کے طور پر جانچنا ہے، کیونکہ $11\times 11=121$، $97$ سے بڑا ہے۔}{$97$ has exactly two \ealkey{factors}, $1$ and itself, so $97$ is a \ealkey{prime number}.\\[0.45em]We only need to test $2$, $3$, $5$ and $7$ as \ealkey{factors}, because $11\times 11=121$ is bigger than $97$.}}

\prq{98}{$1$, $2$, $7$, $14$, $49$, $98$}
  {\ealpara{$98$ کے چھ \ealkeytr{عوامل} ہیں، دو نہیں، اس لیے $98$ \ealkeytr{مفرد عدد} نہیں ہے۔}{$98$ has six \ealkey{factors}, not two, so $98$ is not a \ealkey{prime number}.}}

\prq{99}{$1$, $3$, $9$, $11$, $33$, $99$}
  {\ealpara{$99$ کے چھ \ealkeytr{عوامل} ہیں، دو نہیں، اس لیے $99$ \ealkeytr{مفرد عدد} نہیں ہے۔}{$99$ has six \ealkey{factors}, not two, so $99$ is not a \ealkey{prime number}.}}

\prq{100}{$1$, $2$, $4$, $5$, $10$, $20$, $25$, $50$, $100$}
  {\ealpara{$100$ کے نو \ealkeytr{عوامل} ہیں، دو نہیں، اس لیے $100$ \ealkeytr{مفرد عدد} نہیں ہے۔}{$100$ has nine \ealkey{factors}, not two, so $100$ is not a \ealkey{prime number}.}}

%================================================================
\qtopic{Prime factor decomposition}
%================================================================

\tf[\Large]{\ealpara{$12=2^2\times 3$، $12$ کی \ealkeytr{مفرد عوامل میں تحلیل} ہے}{$12=2^2\times 3$ is the \ealkey{prime factor decomposition} of $12$}}
  {\ealpara{\textbf{سچ۔} $2$ اور $3$ دونوں \ealkeytr{مفرد اعداد} ہیں، اور $4\times 3=12$۔}{\textbf{True.} $2$ and $3$ are both \ealkey{prime}, and $4\times 3=12$.}}

\tf[\Large]{\ealpara{$2\times 3\times 4$، $24$ کی \ealkeytr{مفرد عوامل میں تحلیل} ہے}{$2\times 3\times 4$ is the \ealkey{prime factor decomposition} of $24$}}
  {\ealpara{\textbf{جھوٹ۔} $4$ \ealkeytr{مفرد عدد} نہیں ہے۔ $24=2^3\times 3$}{\textbf{False.} $4$ is not \ealkey{prime}.\\[0.3em] $24=2^3\times 3$}}

\tf[\Large]{\ealpara{$1$ سے بڑا ہر مکمل عدد \ealkeytr{مفرد اعداد} کے \ealkeytr{حاصل ضرب} کے طور پر لکھا جا سکتا ہے}{Every whole number greater than $1$ can be written as a \ealkey{product} of \ealkey{primes}}}
  {\ealpara{\textbf{سچ۔} ایک \ealkeytr{مفرد عدد} خود اپنی تحلیل ہے، اور کوئی بھی دوسرا عدد اس وقت تک تقسیم ہوتا رہتا ہے جب تک صرف \ealkeytr{مفرد اعداد} باقی نہ رہیں۔}{\textbf{True.} A \ealkey{prime} is its own decomposition, and any other number keeps splitting until only \ealkey{primes} are left.}}

\tf[\Large]{\ealpara{ایک عدد کی دو مختلف \ealkeytr{مفرد عوامل میں تحلیل} ہو سکتی ہیں}{A number can have two different \ealkey{prime factor decompositions}}}
  {\ealpara{\textbf{جھوٹ۔} لکھنے کی ترتیب کے علاوہ، ہر عدد کی بالکل ایک تحلیل ہوتی ہے۔}{\textbf{False.} Apart from the order you write them in, every number has exactly one decomposition.}}

\tf[\Large]{\ealpara{$2^3\times 5=30$}{$2^3\times 5=30$}}
  {\ealpara{\textbf{جھوٹ۔} $2^3\times 5=8\times 5=40$۔ $30=2\times 3\times 5$}{\textbf{False.} $2^3\times 5=8\times 5=40$.\\[0.3em] $30=2\times 3\times 5$}}

\qq[\LARGE]{\ealpara{$12$ کو اس کے \ealkeytr{مفرد عوامل} کے \ealkeytr{حاصل ضرب} کے طور پر لکھیں}{Write $12$ as a \ealkey{product} of its \ealkey{prime factors}}}
  {$12=2\times 2\times 3=2^2\times 3$}

\qq[\LARGE]{\ealpara{$18$ کو اس کے \ealkeytr{مفرد عوامل} کے \ealkeytr{حاصل ضرب} کے طور پر لکھیں}{Write $18$ as a \ealkey{product} of its \ealkey{prime factors}}}
  {$18=2\times 3\times 3=2\times 3^2$}

\qq[\LARGE]{\ealpara{$60$ کو اس کے \ealkeytr{مفرد عوامل} کے \ealkeytr{حاصل ضرب} کے طور پر لکھیں}{Write $60$ as a \ealkey{product} of its \ealkey{prime factors}}}
  {$60=2\times 2\times 3\times 5=2^2\times 3\times 5$}

\qq[\LARGE]{\ealpara{$84$ کو اس کے \ealkeytr{مفرد عوامل} کے \ealkeytr{حاصل ضرب} کے طور پر لکھیں}{Write $84$ as a \ealkey{product} of its \ealkey{prime factors}}}
  {$84=2\times 2\times 3\times 7=2^2\times 3\times 7$}

\qq[\LARGE]{\ealpara{$200$ کو اس کے \ealkeytr{مفرد عوامل} کے \ealkeytr{حاصل ضرب} کے طور پر لکھیں}{Write $200$ as a \ealkey{product} of its \ealkey{prime factors}}}
  {$200=2\times 2\times 2\times 5\times 5=2^3\times 5^2$}

\qq[\LARGE]{\ealpara{$360$ کو اس کے \ealkeytr{مفرد عوامل} کے \ealkeytr{حاصل ضرب} کے طور پر لکھیں}{Write $360$ as a \ealkey{product} of its \ealkey{prime factors}}}
  {$360=2^3\times 3^2\times 5$}

\qq[\LARGE]{\ealpara{$1001$ کو اس کے \ealkeytr{مفرد عوامل} کے \ealkeytr{حاصل ضرب} کے طور پر لکھیں}{Write $1001$ as a \ealkey{product} of its \ealkey{prime factors}}}
  {$1001=7\times 143=7\times 11\times 13$}

\qq[\LARGE]{\ealpara{$2^3\times 3\times 5^2$ کو ایک عام عدد کے طور پر حل کریں}{Work out $2^3\times 3\times 5^2$ as an ordinary number}}
  {$8\times 3\times 25=600$}

\qq[\Large][\Large]{\ealpara{$n=2^3\times 3\times 5$۔ $2n$ کو \ealkeytr{مفرد عوامل} کے \ealkeytr{حاصل ضرب} کے طور پر لکھیں۔}{$n=2^3\times 3\times 5$.\\[0.4em] Write $2n$ as a \ealkey{product} of \ealkey{prime factors}.}}
  {\ealpara{دوگنا کرنے سے $2$ کا ایک اور \ealkeytr{عامل} بڑھ جاتا ہے: $2n=2^4\times 3\times 5$}{Doubling adds one more \ealkey{factor} of $2$:\\[0.3em] $2n=2^4\times 3\times 5$}}

\qq[\Large][\Large]{\ealpara{$720=2^4\times 3^2\times 5$۔ $720$ کے کتنے \ealkeytr{عوامل} ہیں؟}{$720=2^4\times 3^2\times 5$.\\[0.4em] How many \ealkey{factors} does $720$ have?}}
  {\ealpara{ہر \ealkeytr{عامل} $0$ سے $4$ تک دو، $0$ سے $2$ تک تین اور $0$ یا $1$ پانچ استعمال کرتا ہے: $5\times 3\times 2=30$ \ealkeytr{عوامل}}{Each \ealkey{factor} uses $0$ to $4$ twos, $0$ to $2$ threes and $0$ or $1$ five:\\[0.3em] $5\times 3\times 2=30$ \ealkey{factors}}}

%================================================================
\qtopic{HCF and LCM from lists}
%================================================================

\hcflcm{4}{6}
  {$1$, \hcfmark{2}, $4$}
  {$1$, \hcfmark{2}, $3$, $6$}
  {$4$, $8$, \lcmmark{12}, $16$, $20$,\\[0.45em]
   $24$, $28$, $32$, $36$, $40$}
  {$6$, \lcmmark{12}, $18$, $24$, $30$,\\[0.45em]
   $36$, $42$, $48$, $54$, $60$}

\hcflcm{6}{8}
  {$1$, \hcfmark{2}, $3$, $6$}
  {$1$, \hcfmark{2}, $4$, $8$}
  {$6$, $12$, $18$, \lcmmark{24}, $30$,\\[0.45em]
   $36$, $42$, $48$, $54$, $60$}
  {$8$, $16$, \lcmmark{24}, $32$, $40$,\\[0.45em]
   $48$, $56$, $64$, $72$, $80$}

\hcflcm{6}{9}
  {$1$, $2$, \hcfmark{3}, $6$}
  {$1$, \hcfmark{3}, $9$}
  {$6$, $12$, \lcmmark{18}, $24$, $30$,\\[0.45em]
   $36$, $42$, $48$, $54$, $60$}
  {$9$, \lcmmark{18}, $27$, $36$, $45$,\\[0.45em]
   $54$, $63$, $72$, $81$, $90$}

\hcflcm{8}{12}
  {$1$, $2$, \hcfmark{4}, $8$}
  {$1$, $2$, $3$, \hcfmark{4}, $6$, $12$}
  {$8$, $16$, \lcmmark{24}, $32$, $40$,\\[0.45em]
   $48$, $56$, $64$, $72$, $80$}
  {$12$, \lcmmark{24}, $36$, $48$, $60$,\\[0.45em]
   $72$, $84$, $96$, $108$, $120$}

\hcflcm{10}{15}
  {$1$, $2$, \hcfmark{5}, $10$}
  {$1$, $3$, \hcfmark{5}, $15$}
  {$10$, $20$, \lcmmark{30}, $40$, $50$,\\[0.45em]
   $60$, $70$, $80$, $90$, $100$}
  {$15$, \lcmmark{30}, $45$, $60$, $75$,\\[0.45em]
   $90$, $105$, $120$, $135$, $150$}

\hcflcm{9}{12}
  {$1$, \hcfmark{3}, $9$}
  {$1$, $2$, \hcfmark{3}, $4$, $6$, $12$}
  {$9$, $18$, $27$, \lcmmark{36}, $45$,\\[0.45em]
   $54$, $63$, $72$, $81$, $90$}
  {$12$, $24$, \lcmmark{36}, $48$, $60$,\\[0.45em]
   $72$, $84$, $96$, $108$, $120$}

\hcflcm{12}{18}
  {$1$, $2$, $3$, $4$, \hcfmark{6}, $12$}
  {$1$, $2$, $3$, \hcfmark{6}, $9$, $18$}
  {$12$, $24$, \lcmmark{36}, $48$, $60$,\\[0.45em]
   $72$, $84$, $96$, $108$, $120$}
  {$18$, \lcmmark{36}, $54$, $72$, $90$,\\[0.45em]
   $108$, $126$, $144$, $162$, $180$}

\hcflcm{14}{21}
  {$1$, $2$, \hcfmark{7}, $14$}
  {$1$, $3$, \hcfmark{7}, $21$}
  {$14$, $28$, \lcmmark{42}, $56$, $70$,\\[0.45em]
   $84$, $98$, $112$, $126$, $140$}
  {$21$, \lcmmark{42}, $63$, $84$, $105$,\\[0.45em]
   $126$, $147$, $168$, $189$, $210$}

\hcflcm{15}{20}
  {$1$, $3$, \hcfmark{5}, $15$}
  {$1$, $2$, $4$, \hcfmark{5}, $10$, $20$}
  {$15$, $30$, $45$, \lcmmark{60}, $75$,\\[0.45em]
   $90$, $105$, $120$, $135$, $150$}
  {$20$, $40$, \lcmmark{60}, $80$, $100$,\\[0.45em]
   $120$, $140$, $160$, $180$, $200$}

\hcflcm{8}{9}
  {\hcfmark{1}, $2$, $4$, $8$}
  {\hcfmark{1}, $3$, $9$}
  {$8$, $16$, $24$, $32$, $40$,\\[0.45em]
   $48$, $56$, $64$, \lcmmark{72}, $80$}
  {$9$, $18$, $27$, $36$, $45$,\\[0.45em]
   $54$, $63$, \lcmmark{72}, $81$, $90$}

%================================================================
\qtopic{HCF and LCM by prime factor decomposition}
%================================================================

\tf[\Large]{\ealpara{دو اعداد کا \ealkeytr{سب سے بڑا مشترک عامل} ہمیشہ دونوں سے چھوٹا ہوتا ہے}{The \ealkey{HCF} of two numbers is always smaller than both of them}}
  {\ealpara{\textbf{جھوٹ۔} $6$ اور $12$ کا \ealkeytr{سب سے بڑا مشترک عامل} $6$ ہے، جو خود ان میں سے ایک عدد ہے۔}{\textbf{False.} The \ealkey{HCF} of $6$ and $12$ is $6$, which is one of the numbers itself.}}

\tf[\Large]{\ealpara{دو اعداد کا \ealkeytr{سب سے چھوٹا مشترک مضاعف} ہمیشہ دونوں کا \ealkeytr{مضاعف} ہوتا ہے}{The \ealkey{LCM} of two numbers is always a \ealkey{multiple} of both of them}}
  {\ealpara{\textbf{سچ۔} یہ بالکل وہی ہے جو ایک مشترک \ealkeytr{مضاعف} ہوتا ہے؛ \ealkeytr{سب سے چھوٹا مشترک مضاعف} سب سے چھوٹا والا ہے۔}{\textbf{True.} That is exactly what a common \ealkey{multiple} is; the \ealkey{LCM} is the smallest one.}}

\tf[\Large]{\ealpara{اگر $a=2^2\times 3$ اور $b=2\times 3^2$، تو $a$ اور $b$ کا \ealkeytr{سب سے بڑا مشترک عامل} $2\times 3$ ہے}{If $a=2^2\times 3$ and $b=2\times 3^2$,\\[0.3em]then the \ealkey{HCF} of $a$ and $b$ is $2\times 3$}}
  {\ealpara{\textbf{سچ۔} ہر \ealkeytr{مفرد عدد} کی \emph{سب سے کم} \ealkeytr{قوت} لیں جو دونوں میں آتی ہے: $2^1\times 3^1=6$۔}{\textbf{True.} Take the \emph{lowest} \ealkey{power} of each \ealkey{prime} that appears in both: $2^1\times 3^1=6$.}}

\tf[\Large]{\ealpara{$\mathrm{HCF}\times\mathrm{LCM}$ دو اعداد کا \ealkeytr{حاصل ضرب} دیتا ہے}{$\mathrm{HCF}\times\mathrm{LCM}$ gives the \ealkey{product} of the two numbers}}
  {\ealpara{\textbf{سچ۔} ہر \ealkeytr{مفرد عدد} ایک بار اپنی سب سے کم \ealkeytr{قوت} پر اور ایک بار سب سے زیادہ پر گنا جاتا ہے، جو دونوں اعداد کو گننے کے برابر ہے۔}{\textbf{True.} Every \ealkey{prime} is counted once at its lowest \ealkey{power} and once at its highest, which is the same as counting both numbers.}}

\tf[\Large]{\ealpara{$6$ اور $8$ کا \ealkeytr{سب سے چھوٹا مشترک مضاعف} $48$ ہے}{The \ealkey{lowest common multiple} of $6$ and $8$ is $48$}}
  {\ealpara{\textbf{جھوٹ۔} $6=2\times 3$ اور $8=2^3$، اس لیے \ealkeytr{LCM} $2^3\times 3=24$ ہے۔}{\textbf{False.} $6=2\times 3$ and $8=2^3$, so the \ealkey{LCM} is $2^3\times 3=24$.}}

\qq[\Large][\Large]{\ealpara{$12=2^2\times 3$ اور $18=2\times 3^2$۔ $12$ اور $18$ کا \ealkeytr{سب سے بڑا مشترک عامل} تلاش کریں۔}{$12=2^2\times 3$ \quad and \quad $18=2\times 3^2$\\[0.4em] Find the \ealkey{HCF} of $12$ and $18$.}}
  {\ealpara{ہر مشترک \ealkeytr{مفرد عدد} کی سب سے کم \ealkeytr{قوت}: $\mathrm{HCF}=2\times 3=6$}{Lowest \ealkey{power} of each shared \ealkey{prime}:\\[0.3em] $\mathrm{HCF}=2\times 3=6$}}

\qq[\Large][\Large]{\ealpara{$12=2^2\times 3$ اور $18=2\times 3^2$۔ $12$ اور $18$ کا \ealkeytr{سب سے چھوٹا مشترک مضاعف} تلاش کریں۔}{$12=2^2\times 3$ \quad and \quad $18=2\times 3^2$\\[0.4em] Find the \ealkey{LCM} of $12$ and $18$.}}
  {\ealpara{ہر \ealkeytr{مفرد عدد} کی سب سے زیادہ \ealkeytr{قوت}: $\mathrm{LCM}=2^2\times 3^2=36$}{Highest \ealkey{power} of each \ealkey{prime}:\\[0.3em] $\mathrm{LCM}=2^2\times 3^2=36$}}

\qq[\Large][\Large]{\ealpara{$24=2^3\times 3$ اور $60=2^2\times 3\times 5$۔ $24$ اور $60$ کا \ealkeytr{سب سے بڑا مشترک عامل} تلاش کریں۔}{$24=2^3\times 3$ \quad and \quad $60=2^2\times 3\times 5$\\[0.4em] Find the \ealkey{HCF} of $24$ and $60$.}}
  {$\mathrm{HCF}=2^2\times 3=12$}

\qq[\Large][\Large]{\ealpara{$24=2^3\times 3$ اور $60=2^2\times 3\times 5$۔ $24$ اور $60$ کا \ealkeytr{سب سے چھوٹا مشترک مضاعف} تلاش کریں۔}{$24=2^3\times 3$ \quad and \quad $60=2^2\times 3\times 5$\\[0.4em] Find the \ealkey{LCM} of $24$ and $60$.}}
  {$\mathrm{LCM}=2^3\times 3\times 5=120$}

\qq[\Large][\large]{\ealpara{$84$ اور $120$ کا \ealkeytr{سب سے بڑا مشترک عامل} اور \ealkeytr{سب سے چھوٹا مشترک مضاعف} تلاش کریں}{Find the \ealkey{HCF} and the \ealkey{LCM} of $84$ and $120$}}
  {\ealpara{$84=2^2\times 3\times 7$ \qquad $120=2^3\times 3\times 5$\\[0.4em]
   $\mathrm{HCF}=2^2\times 3=12$\\[0.3em]
   $\mathrm{LCM}=2^3\times 3\times 5\times 7=840$}{$84=2^2\times 3\times 7$ \qquad $120=2^3\times 3\times 5$\\[0.4em]
   $\mathrm{HCF}=2^2\times 3=12$\\[0.3em]
   $\mathrm{LCM}=2^3\times 3\times 5\times 7=840$}}

\qq[\Large][\large]{\ealpara{$126$ اور $180$ کا \ealkeytr{سب سے بڑا مشترک عامل} اور \ealkeytr{سب سے چھوٹا مشترک مضاعف} تلاش کریں}{Find the \ealkey{HCF} and the \ealkey{LCM} of $126$ and $180$}}
  {\ealpara{$126=2\times 3^2\times 7$ \qquad $180=2^2\times 3^2\times 5$\\[0.4em]
   $\mathrm{HCF}=2\times 3^2=18$\\[0.3em]
   $\mathrm{LCM}=2^2\times 3^2\times 5\times 7=1260$}{$126=2\times 3^2\times 7$ \qquad $180=2^2\times 3^2\times 5$\\[0.4em]
   $\mathrm{HCF}=2\times 3^2=18$\\[0.3em]
   $\mathrm{LCM}=2^2\times 3^2\times 5\times 7=1260$}}

\qq[\Large][\large]{\ealpara{$a=2^3\times 3^2\times 5$ اور $b=2^2\times 3\times 5^2$۔ \ealkeytr{سب سے بڑا مشترک عامل} اور \ealkeytr{سب سے چھوٹا مشترک مضاعف}، \ealkeytr{قوت نما شکل} میں تلاش کریں۔}{$a=2^3\times 3^2\times 5$ \quad and \quad $b=2^2\times 3\times 5^2$\\[0.4em] Find the \ealkey{HCF} and the \ealkey{LCM}, in \ealkey{index form}.}}
  {\ealpara{$\mathrm{HCF}=2^2\times 3\times 5$\\[0.3em]
   $\mathrm{LCM}=2^3\times 3^2\times 5^2$}{$\mathrm{HCF}=2^2\times 3\times 5$\\[0.3em]
   $\mathrm{LCM}=2^3\times 3^2\times 5^2$}}

\qq[\Large][\large]{\ealpara{دو اعداد کا \ealkeytr{سب سے بڑا مشترک عامل} $6$ اور \ealkeytr{سب سے چھوٹا مشترک مضاعف} $72$ ہے۔ ایک عدد $24$ ہے۔ دوسرا تلاش کریں۔}{Two numbers have \ealkey{HCF} $6$ and \ealkey{LCM} $72$.\\[0.4em] One of the numbers is $24$. Find the other.}}
  {\ealpara{$\mathrm{HCF}\times\mathrm{LCM}=6\times 72=432$، جو دونوں اعداد کا \ealkeytr{حاصل ضرب} ہے۔ $432\div 24=18$}{$\mathrm{HCF}\times\mathrm{LCM}=6\times 72=432$, which is the \ealkey{product} of the two numbers.\\[0.4em] $432\div 24=18$}}

\qq[\Large][\large]{\ealpara{$2^5\times 3^2\times 11$ اور $2^3\times 3^4\times 7$ کا \ealkeytr{سب سے بڑا مشترک عامل} اور \ealkeytr{سب سے چھوٹا مشترک مضاعف} تلاش کریں۔ اپنے جوابات \ealkeytr{قوت نما شکل} میں چھوڑیں۔}{Find the \ealkey{HCF} and the \ealkey{LCM} of\\[0.4em] $2^5\times 3^2\times 11$ \quad and \quad $2^3\times 3^4\times 7$\\[0.4em] Leave your answers in \ealkey{index form}.}}
  {\ealpara{$\mathrm{HCF}=2^3\times 3^2$\\[0.3em]
   $\mathrm{LCM}=2^5\times 3^4\times 7\times 11$}{$\mathrm{HCF}=2^3\times 3^2$\\[0.3em]
   $\mathrm{LCM}=2^5\times 3^4\times 7\times 11$}}

\qq[\Large][\large]{\ealpara{دو اعداد کا \ealkeytr{سب سے بڑا مشترک عامل} $12$ اور \ealkeytr{سب سے چھوٹا مشترک مضاعف} $180$ ہے۔ اعداد تلاش کریں۔ کیا ایک سے زیادہ جواب ہو سکتا ہے؟}{Two numbers have \ealkey{HCF} $12$ and \ealkey{LCM} $180$.\\[0.4em] Find the numbers. Is there more than one answer?}}
  {\ealpara{ان کا \ealkeytr{حاصل ضرب} $12\times 180=2160$ ہے، اور دونوں $12$ کے \ealkeytr{مضاعف} ہیں۔ $12$ اور $180$ \quad یا \quad $36$ اور $60$\\[0.3em] دونوں جوڑے کام کرتے ہیں، اس لیے ایک سے زیادہ جواب ہے۔}{Their \ealkey{product} is $12\times 180=2160$, and both are \ealkey{multiples} of $12$.\\[0.4em] $12$ and $180$ \quad or \quad $36$ and $60$\\[0.3em] Both pairs work, so there is more than one answer.}}

%================================================================
\qtopic{Square numbers from prime factors}
%================================================================

\tf[\Large]{\ealpara{ایک عدد مربع ہوتا ہے جب اس کی \ealkeytr{مفرد عوامل میں تحلیل} میں ہر \ealkeytr{قوت} جفت ہو}{A number is square exactly when every \ealkey{power} in its \ealkey{prime factor decomposition} is even}}
  {\ealpara{\textbf{سچ۔} ہر جفت \ealkeytr{قوت} کو آدھا کرنے سے \ealkeytr{جذر} ملتا ہے۔}{\textbf{True.} Halving every even \ealkey{power} gives the \ealkey{square root}.}}

\tf[\Large]{\ealpara{$2^3\times 3^2$ ایک \ealkeytr{مربع عدد} ہے}{$2^3\times 3^2$ is a \ealkey{square number}}}
  {\ealpara{\textbf{جھوٹ۔} $2$ کی \ealkeytr{قوت} طاق ہے، اس لیے عدد کو دو برابر حصوں میں تقسیم نہیں کیا جا سکتا۔}{\textbf{False.} The \ealkey{power} of $2$ is odd, so the number cannot be split into two equal halves.}}

\tf[\Large]{\ealpara{$2^4\times 5^6$ ایک \ealkeytr{مربع عدد} ہے}{$2^4\times 5^6$ is a \ealkey{square number}}}
  {\ealpara{\textbf{سچ۔} $2^4\times 5^6=\left(2^2\times 5^3\right)^2$}{\textbf{True.} $2^4\times 5^6=\left(2^2\times 5^3\right)^2$}}

\tf[\Large]{\ealpara{ہر \ealkeytr{مربع عدد} کے \ealkeytr{عوامل} کی تعداد جفت ہوتی ہے}{Every \ealkey{square number} has an even number of \ealkey{factors}}}
  {\ealpara{\textbf{جھوٹ۔} \ealkeytr{مربع اعداد} کے \ealkeytr{عوامل} کی تعداد \emph{طاق} ہوتی ہے: $16$ کے پاس $1,\ 2,\ 4,\ 8,\ 16$ ہیں۔}{\textbf{False.} \ealkey{Square numbers} have an \emph{odd} number of \ealkey{factors}:\\[0.3em] $16$ has $1,\ 2,\ 4,\ 8,\ 16$.}}

\tf[\Large]{\ealpara{اگر $n$ ایک \ealkeytr{مربع عدد} ہے، تو $4n$ بھی ایک \ealkeytr{مربع عدد} ہے}{If $n$ is a \ealkey{square number}, then $4n$ is a \ealkey{square number} too}}
  {\ealpara{\textbf{سچ۔} $4=2^2$ سے ضرب دینے سے ہر \ealkeytr{قوت} جفت رہتی ہے۔}{\textbf{True.} Multiplying by $4=2^2$ keeps every \ealkey{power} even.}}

\qq[\Large][\Large]{\ealpara{\ealkeytr{مفرد عوامل} استعمال کرتے ہوئے دکھائیں کہ $36$ ایک \ealkeytr{مربع عدد} ہے}{Use \ealkey{prime factors} to show that $36$ is a \ealkey{square number}}}
  {\ealpara{$36=2^2\times 3^2$\\[0.3em] دونوں \ealkeytr{قوتیں} جفت ہیں، اور $36=(2\times 3)^2=6^2$۔}{$36=2^2\times 3^2$\\[0.3em] Both \ealkey{powers} are even, and $36=(2\times 3)^2=6^2$.}}

\qq[\Large][\Large]{\ealpara{کیا $2^2\times 3^2\times 5^2$ ایک \ealkeytr{مربع عدد} ہے؟ اگر ہاں، تو اس کا \ealkeytr{جذر} کیا ہے؟}{Is $2^2\times 3^2\times 5^2$ a \ealkey{square number}?\\[0.4em] If so, what is its \ealkey{square root}?}}
  {\ealpara{جی ہاں --- ہر \ealkeytr{قوت} جفت ہے۔ $\sqrt{2^2\times 3^2\times 5^2}=2\times 3\times 5=30$}{Yes --- every \ealkey{power} is even.\\[0.3em] $\sqrt{2^2\times 3^2\times 5^2}=2\times 3\times 5=30$}}

\qq[\Large][\Large]{\ealpara{کیا $2^3\times 3^2$ ایک \ealkeytr{مربع عدد} ہے؟}{Is $2^3\times 3^2$ a \ealkey{square number}?}}
  {\ealpara{نہیں۔ \ealkeytr{قوتوں} کو آدھا کرنے کے لیے $2^{1.5}$ درکار ہوگا، اس لیے $2$ کی \ealkeytr{قوت} جفت ہونی چاہیے اور یہ نہیں ہے۔}{No. Halving the \ealkey{powers} would need $2^{1.5}$,\\[0.3em] so the \ealkey{power} of $2$ must be even and it is not.}}

\qq[\Large][\Large]{\ealpara{$\sqrt{2^6\times 3^4}$ حل کریں}{Work out $\sqrt{2^6\times 3^4}$}}
  {\ealpara{ہر \ealkeytr{قوت} کو آدھا کریں: $2^3\times 3^2=8\times 9=72$}{Halve each \ealkey{power}:\\[0.3em] $2^3\times 3^2=8\times 9=72$}}

\qq[\Large][\Large]{\ealpara{\ealkeytr{مفرد عوامل} استعمال کرتے ہوئے فیصلہ کریں کہ آیا $400$ ایک \ealkeytr{مربع عدد} ہے}{Use \ealkey{prime factors} to decide whether $400$ is a \ealkey{square number}}}
  {\ealpara{$400=2^4\times 5^2$\\[0.3em] دونوں \ealkeytr{قوتیں} جفت ہیں، اس لیے $400=\left(2^2\times 5\right)^2=20^2$۔}{$400=2^4\times 5^2$\\[0.3em] Both \ealkey{powers} are even, so $400=\left(2^2\times 5\right)^2=20^2$.}}

\qq[\Large][\Large]{\ealpara{\ealkeytr{مفرد عوامل} استعمال کرتے ہوئے فیصلہ کریں کہ آیا $150$ ایک \ealkeytr{مربع عدد} ہے}{Use \ealkey{prime factors} to decide whether $150$ is a \ealkey{square number}}}
  {\ealpara{$150=2\times 3\times 5^2$\\[0.3em] $2$ اور $3$ کی \ealkeytr{قوتیں} طاق ہیں، اس لیے $150$ مربع نہیں ہے۔}{$150=2\times 3\times 5^2$\\[0.3em] The \ealkey{powers} of $2$ and $3$ are odd, so $150$ is not square.}}

\qq[\Large][\Large]{\ealpara{$\sqrt{2^8\times 3^2\times 5^4}$ حل کریں}{Work out $\sqrt{2^8\times 3^2\times 5^4}$}}
  {\ealpara{ہر \ealkeytr{قوت} کو آدھا کریں: $2^4\times 3\times 5^2=16\times 3\times 25=1200$}{Halve each \ealkey{power}:\\[0.3em] $2^4\times 3\times 5^2=16\times 3\times 25=1200$}}

\qq[\Large][\large]{\ealpara{$n=2^3\times 3\times 5^2$۔ سب سے چھوٹا مکمل عدد $k$ تلاش کریں تاکہ $nk$ ایک \ealkeytr{مربع عدد} ہو۔}{$n=2^3\times 3\times 5^2$\\[0.4em] Find the smallest whole number $k$ so that $nk$ is a \ealkey{square number}.}}
  {\ealpara{$2$ اور $3$ کی \ealkeytr{قوتیں} طاق ہیں، اس لیے ہر ایک کو ایک اور چاہیے: $k=2\times 3=6$، جس سے $nk=2^4\times 3^2\times 5^2$۔}{The \ealkey{powers} of $2$ and $3$ are odd, so each needs one more:\\[0.3em] $k=2\times 3=6$, giving $nk=2^4\times 3^2\times 5^2$.}}

\qq[\Large][\large]{\ealpara{$504=2^3\times 3^2\times 7$۔ سب سے چھوٹا عدد تلاش کریں جس سے $504$ کو ضرب دے کر \ealkeytr{مربع عدد} حاصل ہو۔}{$504=2^3\times 3^2\times 7$\\[0.4em] Find the smallest number $504$ can be multiplied by to give a \ealkey{square number}.}}
  {\ealpara{$2$ اور $7$ کی \ealkeytr{قوتیں} طاق ہیں: $2\times 7=14$ سے ضرب دیں۔ $504\times 14=7056=84^2$}{The \ealkey{powers} of $2$ and $7$ are odd:\\[0.3em] multiply by $2\times 7=14$.\\[0.3em] $504\times 14=7056=84^2$}}

\qq[\Large][\large]{\ealpara{$1080=2^3\times 3^3\times 5$۔ سب سے چھوٹا عدد تلاش کریں جس سے $1080$ کو تقسیم کر کے \ealkeytr{مربع عدد} حاصل ہو۔}{$1080=2^3\times 3^3\times 5$\\[0.4em] Find the smallest number $1080$ can be divided by to give a \ealkey{square number}.}}
  {\ealpara{ایک $2$، ایک $3$ اور $5$ ہٹا دیں: $2\times 3\times 5=30$ سے تقسیم کریں۔ $1080\div 30=36=6^2$}{Remove one $2$, one $3$ and the $5$:\\[0.3em] divide by $2\times 3\times 5=30$.\\[0.3em] $1080\div 30=36=6^2$}}

%================================================================
\qtopic{Problem solving with HCF and LCM}
%================================================================

\tf[\Large]{\ealpara{یہ جاننے کے لیے کہ دو دہرائے جانے والے واقعات اگلی بار ایک ساتھ کب ہوں گے، آپ \ealkeytr{سب سے چھوٹا مشترک مضاعف} استعمال کرتے ہیں}{To find when two repeating events next happen together, you use the \ealkey{LCM}}}
  {\ealpara{\textbf{سچ۔} آپ پہلا وقت تلاش کر رہے ہیں جو دونوں وقفوں کا \ealkeytr{مضاعف} ہو۔}{\textbf{True.} You are looking for the first time that is a \ealkey{multiple} of both gaps.}}

\tf[\Large]{\ealpara{دو مجموعوں کو سب سے بڑے ایک جیسے تھیلوں میں تقسیم کرنے کے لیے، آپ \ealkeytr{سب سے چھوٹا مشترک مضاعف} استعمال کرتے ہیں}{To split two collections into the largest identical bags, you use the \ealkey{LCM}}}
  {\ealpara{\textbf{جھوٹ۔} \ealkeytr{سب سے بڑا مشترک عامل} استعمال کریں --- تھیلے کا حجم دونوں مقداروں کا \ealkeytr{عامل} ہونا چاہیے۔}{\textbf{False.} Use the \ealkey{HCF} --- the bag size has to be a \ealkey{factor} of both amounts.}}

\tf[\Large]{\ealpara{دو رسیوں کو برابر ٹکڑوں میں کاٹنا، جتنا لمبا ممکن ہو اور کچھ بچے نہ، \ealkeytr{سب سے بڑا مشترک عامل} استعمال کرتا ہے}{Cutting two ropes into equal pieces, as long as possible with none left over, uses the \ealkey{HCF}}}
  {\ealpara{\textbf{سچ۔} ٹکڑے کی لمبائی دونوں رسیوں کی لمبائی کا \ealkeytr{عامل} ہونی چاہیے۔}{\textbf{True.} The piece length must be a \ealkey{factor} of both rope lengths.}}

\tf[\Large]{\ealpara{دو اعداد کا \ealkeytr{سب سے چھوٹا مشترک مضاعف} کبھی بھی ان میں سے کسی سے چھوٹا نہیں ہوتا}{The \ealkey{LCM} of two numbers is never smaller than either of them}}
  {\ealpara{\textbf{سچ۔} یہ ہر عدد کا \ealkeytr{مضاعف} ہے، اس لیے یہ کم از کم دونوں کے برابر بڑا ہوتا ہے۔}{\textbf{True.} It is a \ealkey{multiple} of each number, so it is at least as big as both.}}

\tf[\Large]{\ealpara{دو واقعات ہر $4$ منٹ اور ہر $6$ منٹ بعد دہرائے جاتے ہیں، اس لیے وہ ہر $24$ منٹ بعد ملتے ہیں}{Two events repeat every $4$ minutes and every $6$ minutes,\\[0.3em] so they coincide every $24$ minutes}}
  {\ealpara{\textbf{جھوٹ۔} وہ ہر $\mathrm{LCM}(4,6)=12$ منٹ بعد ملتے ہیں۔}{\textbf{False.} They coincide every $\mathrm{LCM}(4,6)=12$ minutes.}}

\qq[\Large][\Large]{\ealpara{دو گھنٹیاں اب ایک ساتھ بجتی ہیں۔ ایک ہر $4$ منٹ میں، دوسری ہر $6$ منٹ میں۔ وہ اگلی بار ایک ساتھ کب بجیں گی؟}{Two bells ring together now.\\[0.4em] One rings every $4$ minutes, the other every $6$ minutes.\\[0.4em] When do they next ring together?}}
  {\ealpara{مشترک \ealkeytr{مضاعف}، اس لیے \ealkeytr{LCM} استعمال کریں: $\mathrm{LCM}(4,6)=12$ --- $12$ منٹ میں۔}{Common \ealkey{multiple}, so use the \ealkey{LCM}:\\[0.3em] $\mathrm{LCM}(4,6)=12$ --- in $12$ minutes.}}

\qq[\Large][\Large]{\ealpara{دو ربن، $18$\,cm اور $24$\,cm لمبے، برابر ٹکڑوں میں کاٹے جاتے ہیں، جتنا لمبا ممکن ہو اور کچھ بچے نہ۔ ہر ٹکڑا کتنا لمبا ہے؟}{Two ribbons, $18$\,cm and $24$\,cm long, are cut into equal pieces, as long as possible with none left over.\\[0.4em] How long is each piece?}}
  {\ealpara{لمبائی دونوں کی \ealkeytr{عامل} ہونی چاہیے، اس لیے \ealkeytr{HCF} استعمال کریں: $\mathrm{HCF}(18,24)=6$ --- ہر ٹکڑا $6$\,cm ہے۔}{The length must be a \ealkey{factor} of both, so use the \ealkey{HCF}:\\[0.3em] $\mathrm{HCF}(18,24)=6$ --- each piece is $6$\,cm.}}

\qq[\Large][\large]{\ealpara{$36$ پین اور $48$ پنسلیں ایک جیسے پیکٹوں میں بغیر کچھ بچے کے تقسیم کی جاتی ہیں۔ پیکٹوں کی سب سے بڑی تعداد کیا ہے؟}{$36$ pens and $48$ pencils are shared into identical packs with nothing left over.\\[0.4em] What is the greatest number of packs?}}
  {\ealpara{$36=2^2\times 3^2$ \qquad $48=2^4\times 3$\\[0.4em] $\mathrm{HCF}=2^2\times 3=12$ پیکٹ، ہر ایک میں $3$ پین اور $4$ پنسلیں۔}{$36=2^2\times 3^2$ \qquad $48=2^4\times 3$\\[0.4em] $\mathrm{HCF}=2^2\times 3=12$ packs,\\[0.3em] each with $3$ pens and $4$ pencils.}}

\qq[\Large][\Large]{\ealpara{دو لائٹ ہاؤس $8$pm پر ایک ساتھ چمکتے ہیں۔ ایک ہر $15$ سیکنڈ میں، دوسرا ہر $18$ سیکنڈ میں۔ وہ اگلی بار ایک ساتھ کب چمکیں گے؟}{Two lighthouses flash together at $8$pm.\\[0.4em] One flashes every $15$ seconds, the other every $18$ seconds.\\[0.4em] When do they next flash together?}}
  {\ealpara{$\mathrm{LCM}(15,18)=90$ سیکنڈ\\[0.3em] $8${:}$01${:}$30$pm پر۔}{$\mathrm{LCM}(15,18)=90$ seconds\\[0.3em] At $8${:}$01${:}$30$pm.}}

\qq[\Large][\Large]{\ealpara{مستطیل ٹائلیں $12$\,cm بائی $18$\,cm ہیں۔ سب سے چھوٹا مربع کیا ہے جسے وہ بالکل ٹائل کر سکیں؟}{Rectangular tiles measure $12$\,cm by $18$\,cm.\\[0.4em] What is the smallest square they can tile exactly?}}
  {\ealpara{ضلع دونوں کا \ealkeytr{مضاعف} ہونا چاہیے، اس لیے \ealkeytr{LCM} استعمال کریں: $\mathrm{LCM}(12,18)=36$\\[0.3em] $36$\,cm بائی $36$\,cm مربع۔}{The side must be a \ealkey{multiple} of both, so use the \ealkey{LCM}:\\[0.3em] $\mathrm{LCM}(12,18)=36$\\[0.3em] A $36$\,cm by $36$\,cm square.}}

\qq[\Large][\Large]{\ealpara{ایک بس ہر $24$ منٹ میں اور ایک ٹرام ہر $36$ منٹ میں روانہ ہوتی ہے۔ دونوں $9${:}$00$ پر روانہ ہوتی ہیں۔ وہ اگلی بار ایک ساتھ کب روانہ ہوں گی؟}{A bus leaves every $24$ minutes and a tram every $36$ minutes.\\[0.4em] Both leave at $9${:}$00$. When do they next leave together?}}
  {\ealpara{$\mathrm{LCM}(24,36)=72$ منٹ\\[0.3em] $10${:}$12$ پر۔}{$\mathrm{LCM}(24,36)=72$ minutes\\[0.3em] At $10${:}$12$.}}

\qq[\Large][\large]{\ealpara{دو کوگ، $18$ اور $24$ دانتوں والے، ایک نشان زدہ دانت پر شروع میں جڑے ہوئے ہیں۔ نشان دوبارہ ملنے سے پہلے ہر ایک کتنے چکر لگاتا ہے؟}{Two cogs, with $18$ and $24$ teeth, start meshed at a marked tooth.\\[0.4em] How many turns does each make before the marks meet again?}}
  {\ealpara{$\mathrm{LCM}(18,24)=72$ دانت گزرنے چاہئیں۔ $72\div 18=4$ چکر چھوٹے کوگ کے، $72\div 24=3$ چکر بڑے کوگ کے۔}{$\mathrm{LCM}(18,24)=72$ teeth must pass.\\[0.4em] $72\div 18=4$ turns of the small cog,\\[0.3em] $72\div 24=3$ turns of the large cog.}}

\qq[\Large][\large]{\ealpara{ایک فرش $360$\,cm بائی $504$\,cm ہے۔ اسے ایک جیسی مربع ٹائلوں سے ٹائل کیا جاتا ہے، جتنی بڑی ممکن ہو۔ ٹائل کا سائز اور ٹائلوں کی تعداد تلاش کریں۔}{A floor measures $360$\,cm by $504$\,cm.\\[0.4em] It is tiled with identical square tiles, as large as possible.\\[0.4em] Find the tile size and the number of tiles.}}
  {\ealpara{$360=2^3\times 3^2\times 5$ \qquad $504=2^3\times 3^2\times 7$\\[0.4em] $\mathrm{HCF}=2^3\times 3^2=72$، اس لیے ٹائلیں $72$\,cm مربع ہیں۔ $5\times 7=35$ ٹائلیں۔}{$360=2^3\times 3^2\times 5$ \qquad $504=2^3\times 3^2\times 7$\\[0.4em] $\mathrm{HCF}=2^3\times 3^2=72$, so the tiles are $72$\,cm square.\\[0.3em] $5\times 7=35$ tiles.}}

\qq[\Large][\large]{\ealpara{تین دوڑنے والے $60$\,s، $72$\,s اور $90$\,s فی چکر لیتے ہیں۔ وہ ایک ساتھ شروع کرتے ہیں۔ وہ اگلی بار سب کے سب اسٹارٹ لائن پر کب ہوں گے؟}{Three runners take $60$\,s, $72$\,s and $90$\,s per lap.\\[0.4em] They start together. When are they next all at the start line?}}
  {\ealpara{$60=2^2\times 3\times 5$ \quad $72=2^3\times 3^2$ \quad $90=2\times 3^2\times 5$\\[0.4em] $\mathrm{LCM}=2^3\times 3^2\times 5=360$\,s\\[0.3em] $6$ منٹ بعد۔}{$60=2^2\times 3\times 5$ \quad $72=2^3\times 3^2$ \quad $90=2\times 3^2\times 5$\\[0.4em] $\mathrm{LCM}=2^3\times 3^2\times 5=360$\,s\\[0.3em] After $6$ minutes.}}

\qq[\Large][\large]{\ealpara{ہاٹ ڈاگ $10$ کے پیکٹوں میں، بن $8$ کے پیکٹوں میں اور ساس ساشے $12$ کے پیکٹوں میں آتے ہیں۔ سام کے لیے بغیر کچھ بچے کے بنائے جا سکنے والے مکمل ہاٹ ڈاگ کی سب سے چھوٹی تعداد کیا ہے، اور یہ ہر ایک کے کتنے پیکٹ ہیں؟}{Hot dogs come in packs of $10$, buns in packs of $8$ and sauce sachets in packs of $12$.\\[0.4em] What is the smallest number of complete hot dogs Sam can make with nothing left over, and how many packs of each is that?}}
  {\ealpara{$\mathrm{LCM}(10,8,12)=2^3\times 3\times 5=120$\\[0.4em] $120$ ہاٹ ڈاگ: $12$ پیکٹ ہاٹ ڈاگ، $15$ پیکٹ بن اور $10$ پیکٹ ساشے۔}{$\mathrm{LCM}(10,8,12)=2^3\times 3\times 5=120$\\[0.4em] $120$ hot dogs: $12$ packs of hot dogs,\\[0.3em] $15$ packs of buns and $10$ packs of sachets.}}

\end{document}